\documentclass[twoside,leqno]{article}

% --------------------------------------------------------------------------
% USER-ADJUSTABLE SETTINGS
% Change the values in this one block to retune the page and type area.  The
% subordinate font sizes and their leading scale automatically with the body
% size; no page-specific breaks or fixed vertical fillers are used below.
\providecommand{\DocumentPaperWidth}{176mm}  % ISO B5
\providecommand{\DocumentPaperHeight}{250mm}
\providecommand{\DocumentLeftMargin}{0.7cm}
\providecommand{\DocumentRightMargin}{0.7cm}
\providecommand{\DocumentTopMargin}{0.91cm}  % includes the running page number
\providecommand{\DocumentBottomMargin}{0.91cm}
\providecommand{\DocumentHeaderHeight}{16pt}
\providecommand{\DocumentHeaderSeparation}{4pt}
\providecommand{\DocumentBodyPointSize}{13}  % numerical value, in pt
\providecommand{\DocumentLineSpread}{1.1}
\providecommand{\DocumentEmergencyStretch}{5em} % flexible line-breaking reserve
\providecommand{\DocumentChapterBeforeSkip}{0.6\baselineskip}
\providecommand{\DocumentChapterAfterSkip}{\smallskipamount}
\providecommand{\DocumentSectionBeforeSkip}{0.4\baselineskip}
\providecommand{\DocumentTitleToBySkip}{0.6ex}
\providecommand{\DocumentByToAuthorSkip}{0.2ex}
\providecommand{\DocumentTitleBlockAfterSkip}{0.7\baselineskip}
% --------------------------------------------------------------------------

\usepackage{xfp}
\usepackage{geometry}
\geometry{
  paperwidth=\DocumentPaperWidth,
  paperheight=\DocumentPaperHeight,
  left=\DocumentLeftMargin,
  right=\DocumentRightMargin,
  top=\DocumentTopMargin,
  bottom=\DocumentBottomMargin,
  includehead,
  headheight=\DocumentHeaderHeight,
  headsep=\DocumentHeaderSeparation
}
\usepackage[T1]{fontenc}
\usepackage{newpxtext}
\usepackage{newpxmath}
\usepackage{amsmath,mathtools}
\usepackage{bm}
\usepackage{microtype}
\usepackage{xcolor}
\usepackage{fancyhdr}
\usepackage{enumitem}
\usepackage{xurl}
\usepackage[unicode,pdfencoding=auto,psdextra]{hyperref}

\definecolor{LinkBlue}{HTML}{1C4E80}
\definecolor{CiteViolet}{HTML}{6A3D9A}
\definecolor{UrlTeal}{HTML}{00796B}
\hypersetup{
  colorlinks=true,
  linkcolor=LinkBlue,
  citecolor=CiteViolet,
  urlcolor=UrlTeal,
  pdfborder={0 0 0},
  pdftitle={On Rings of Operators. II},
  pdfauthor={F. J. Murray and J. von Neumann},
  pdfsubject={Faithful LaTeX transcription of the 1937 paper},
  pdfdisplaydoctitle=true,
  bookmarksopen=true,
  bookmarksopenlevel=2
}
\urlstyle{same}

% A genuine body face at \DocumentBodyPointSize, with proportionally scaled
% auxiliary sizes.  newpxtext and newpxmath should be changed as a pair if a
% user substitutes another text/math family.
\newcommand{\ScaledFontPoint}[1]{\fpeval{\DocumentBodyPointSize*(#1)}pt}
\makeatletter
\renewcommand\normalsize{%
  \@setfontsize\normalsize{\ScaledFontPoint{1}}{\ScaledFontPoint{6/5}}%
  \abovedisplayskip \ScaledFontPoint{10/13} plus \ScaledFontPoint{2/13} minus \ScaledFontPoint{5/13}%
  \abovedisplayshortskip 0pt plus \ScaledFontPoint{3/13}%
  \belowdisplayshortskip \ScaledFontPoint{6/13} plus \ScaledFontPoint{3/13} minus \ScaledFontPoint{3/13}%
  \belowdisplayskip \abovedisplayskip
  \let\@listi\@listI}
\renewcommand\small{\@setfontsize\small{\ScaledFontPoint{23/26}}{\ScaledFontPoint{69/65}}}
\renewcommand\footnotesize{\@setfontsize\footnotesize{\ScaledFontPoint{10/13}}{\ScaledFontPoint{12/13}}}
\renewcommand\scriptsize{\@setfontsize\scriptsize{\ScaledFontPoint{9/13}}{\ScaledFontPoint{54/65}}}
\renewcommand\tiny{\@setfontsize\tiny{\ScaledFontPoint{8/13}}{\ScaledFontPoint{48/65}}}
\renewcommand\large{\@setfontsize\large{\ScaledFontPoint{15/13}}{\ScaledFontPoint{18/13}}}
\renewcommand\Large{\@setfontsize\Large{\ScaledFontPoint{17/13}}{\ScaledFontPoint{102/65}}}
\renewcommand\LARGE{\@setfontsize\LARGE{\ScaledFontPoint{20/13}}{\ScaledFontPoint{24/13}}}
\makeatother
\linespread{\DocumentLineSpread}
\AtBeginDocument{\normalsize}

\setlength{\parindent}{1.45em}
\setlength{\parskip}{0pt}
\setlist{nosep,leftmargin=2em}
\setcounter{secnumdepth}{0}
\emergencystretch=\DocumentEmergencyStretch

% Long aligned displays may break when genuinely needed; short displays normally move intact.
\allowdisplaybreaks[1]
% Keep vertical spacing intrinsic instead of stretching it to fill each page.
\raggedbottom

% Outer-corner page numbers; the first page is suppressed below.
\pagestyle{fancy}
\fancyhf{}
\fancyhead[LE,RO]{\footnotesize\thepage}
\renewcommand{\headrulewidth}{0pt}
\renewcommand{\footrulewidth}{0pt}

% Semantic helpers used throughout the transcription.
% The source sets both letters of Tr in mathematical italic, not as an upright
% operator name; keep that distinction visible throughout.
\newcommand{\Tr}{Tr}
\newcommand{\TrM}{Tr_M}
\newcommand{\TrMp}{Tr_{M'}}
\newcommand{\DM}{D_{M}}
\newcommand{\DMp}{D_{M'}}
\newcommand{\cH}{\mathfrak{H}}
\newcommand{\Q}{\mathcal{Q}}
% The source distinguishes this three-stroke operator bound from the ordinary
% two-stroke vector norm.  Negative kerns make the strokes merge in newpx.
\newcommand{\opnorm}[1]{\lvert\lvert\lvert#1\rvert\rvert\rvert}
\newcommand{\RO}{\href{https://doi.org/10.2307/1968693}{R.O.}}
\newcommand{\RORef}[1]{\href{https://doi.org/10.2307/1968693}{#1}}
\makeatletter
\expandafter\def\csname ROlink@4\endcsname#1{\href{https://doi.org/10.1007/978-3-663-15768-7}{#1}}
\expandafter\def\csname ROlink@15\endcsname#1{\href{https://doi.org/10.1090/coll/015}{#1}}
\expandafter\def\csname ROlink@16\endcsname#1{\href{https://doi.org/10.1007/BF01782338}{#1}}
\expandafter\def\csname ROlink@18\endcsname#1{\href{https://doi.org/10.1007/BF01782352}{#1}}
\expandafter\def\csname ROlink@19\endcsname#1{\href{https://gdz.sub.uni-goettingen.de/download/pdf/PPN379400774/PPN379400774.pdf}{#1}}
\expandafter\def\csname ROlink@22\endcsname#1{\href{https://doi.org/10.2307/1968692}{#1}}
\newcommand{\ROcite}[1]{%
  \@ifundefined{ROlink@#1}
    {\href{https://doi.org/10.2307/1968693}{(#1)}}
    {\csname ROlink@#1\endcsname{(#1)}}% 
}
\makeatother
\newcommand{\SectionTarget}[1]{\phantomsection\hypertarget{sec:#1}{}}
\newcommand{\SectionStart}[2]{%
  \par\addvspace{\DocumentSectionBeforeSkip}\noindent
  \SectionTarget{#1}\pdfbookmark[2]{#2}{outline:sec:#1}%
  \hspace*{\parindent}\ignorespaces
}
\newcommand{\SectionRef}[1]{\hyperlink{sec:#1}{\S#1}}
\newcommand{\ChapterTarget}[1]{\phantomsection\hypertarget{chap:#1}{}}
\newcommand{\ChapterRef}[1]{\hyperlink{chap:#1}{Chapter~#1}}
\newcommand{\StatementTarget}[1]{\phantomsection\hypertarget{stmt:#1}{}}
\newcommand{\StatementRef}[2]{\hyperlink{stmt:#1}{#2}}
\newcommand{\EquationTarget}[1]{\hypertarget{eq:#1}{}}
\newcommand{\EquationRef}[2]{\hyperlink{eq:#1}{#2}}
\newenvironment{statement}[2]{%
  \par\medskip
  \noindent\StatementTarget{#1}\hspace*{\parindent}\textbf{#2.}\ \itshape\ignorespaces
}{\par}
\newcommand{\proofhead}{\par\smallskip\noindent\textit{Proof.}\ \ignorespaces}
\newcommand{\chapterhead}[3]{%
  \par\addvspace{\DocumentChapterBeforeSkip}\noindent\ChapterTarget{#1}%
  \pdfbookmark[1]{#3}{outline:chap:#1}%
  \makebox[\linewidth][c]{\bfseries\MakeUppercase{Chapter #1. #2}}%
  \par\addvspace{\DocumentChapterAfterSkip}
}
\newcommand{\appendixhead}{%
  \par\addvspace{\DocumentChapterBeforeSkip}\noindent\ChapterTarget{Appendix}%
  \pdfbookmark[1]{Appendix. The matrix-aspect of the M of class (II₁)}{outline:chap:Appendix}%
  \makebox[\linewidth][c]{\bfseries APPENDIX}%
  \par\addvspace{\DocumentChapterAfterSkip}
}

\begin{document}
\setcounter{page}{1}
\thispagestyle{empty}
\hypertarget{doc-start}{}
\pdfbookmark[0]{On Rings of Operators. II}{outline:document}

\begingroup
\renewcommand{\thefootnote}{\fnsymbol{footnote}}
\par\noindent
{\centering
  {\Large\bfseries ON RINGS OF OPERATORS. II\footnotemark[1]}\\[\DocumentTitleToBySkip]
  {\small BY}\\[\DocumentByToAuthorSkip]
  {\bfseries F. J. MURRAY\footnotemark[2] AND J. von NEUMANN}\par}
\footnotetext[1]{Presented to the Society, October 31, 1936; received by the editors July 28, 1936.}
\footnotetext[2]{National Research Fellow.}
\endgroup
\vspace{\DocumentTitleBlockAfterSkip}

\par\noindent\hypertarget{sec:introduction}{}%
\pdfbookmark[1]{Introduction}{outline:introduction}%
\hspace*{\parindent}%
\textbf{Introduction.} This paper is a continuation of one by the same authors:
\href{https://doi.org/10.2307/1968693}{\textit{On rings of operators}},
Annals of Mathematics, (2), vol.~37 (1936), pp.~116--229.
% Editorial note: source prints the apparently superfluous comma in which, were; retained verbatim.
It contains the solution of certain problems which, were left open there. We
will prove the general additivity of trace $\TrM(A)$, its weak continuity, and
certain isomorphisms between $\cH$, $M$, and $M'$ (cf. the remarks (i)--(iv) at the
end of the above quoted paper). All these considerations refer to ``Case (II)''
for $M$ (cf. \RORef{Theorem VIII, loc. cit.}).

The properties of $\TrM(A)$ are established by obtaining for it a representation
\[
  \TrM(A)=\sum_{i=1}^{m}(Ag_i,g_i)
\]
(with a fixed, finite $m=1,2,\cdots$, and fixed $g_1,\cdots,g_m\in\cH$). This representation
is remarkable, because it is obviously a close analogue of the representation
of $\TrM(A)$ as a trace, that is, as the arithmetic mean of the diagonal
matrix-elements of $A$ in the cases $(\mathrm I_n)$, $n=1,2,\cdots$, when $M$ is essentially
% Editorial note: the unusual source wording ``n-(finite-) dimensional'' is retained verbatim.
the full matrix ring of an $n$-(finite-) dimensional Euclidean space.

For certain cases (with the help of which the others are then mastered)
we have even $m=1$.

In \hyperlink{chap:I}{Part I} the above representation of $\TrM(A)$ is obtained approximately.
The technically interested reader may find it worth observing that the exhaustion
method we use there (\hyperlink{sec:1.2}{\S\S1.2} and \hyperlink{sec:1.3}{1.3}) is analogous to certain procedures
which can be used advantageously in the theories of measures and integration too.
On this basis we establish the main properties of $\TrM(A)$ in \hyperlink{chap:II}{Part II},
and then obtain the exact representation of $\TrM(A)$ in \hyperlink{chap:III}{Part III}.
Here two maximum-problems, called (A) and (B), which seem to possess some
independent interest too, play a decisive role.

\hyperlink{chap:IV}{Part IV} is devoted to establishing an isomorphism between $\cH$, $M$, and $M'$.
It turns out that a certain algebraic-topological extension $\Q(M)$ of $M$ is
isomorphic to $\cH$ and that $M$ and $M'$ play in it the role of right- and
left-multiplication. This leads to an interesting and entirely new type of infinite
hypercomplex systems, which are at the same time Hilbert spaces. A subsequent
paper will be devoted to their independent study.

The appendix deals with the possibility of considering $M$ (in the case
$(\mathrm{II}_1)$) as a system of matrices with continuously spread rows and columns.

We will use the notations, definitions, and results of our paper
\href{https://doi.org/10.2307/1968693}{\textit{On rings of operators}}, quoted
above, throughout this paper. We will quote it, whenever
necessary, as \RO{} All other quotations follow the bibliography of \RO{}
(pp.~125--126, Nos.~(1)--(22)).

The isomorphism problems of different rings $M$ of class $(\mathrm{II}_1)$ (cf. the
remark (v) at the end of \RO{}) are not discussed here. They will be dealt with
in a subsequent publication.

Since the appearance of \RO{} the second-named author has succeeded in
finding new representations of case $(\mathrm{II}_1)$ in terms of infinite direct products,
which throw new light on $(\mathrm{II}_1)$ as a limiting case of the $(\mathrm I_n)$, $n=1,2,\cdots$,
as well as a way of applying the present theory to quantum-mechanics. These
subjects too will be discussed in papers which will follow soon.

\begin{center}
  \hypertarget{contents}{}\pdfbookmark[1]{Contents}{outline:contents}\textsc{Contents}
\end{center}
\noindent\hyperlink{sec:introduction}{Introduction.}\\
\hyperlink{chap:I}{Chapter I. Approximate form of $\TrM(A)$ (for $\alpha\geq1$).}\\
\hspace*{2em}\hyperlink{sec:1.1}{1.1. The normalization.}\\
\hspace*{2em}\hyperlink{sec:1.2}{1.2. The notions $E\geq\lambda$, $E\leq\lambda$ and
$E\mathrel{\underset p\geq}\lambda$, $E\mathrel{\underset p\leq}\lambda$.}\\
\hspace*{2em}\hyperlink{sec:1.3}{1.3. Selection of an approximately homogeneous piece.}\\
\hspace*{2em}\hyperlink{sec:1.4}{1.4. Properties of such a piece.}\\
\hspace*{2em}\hyperlink{sec:1.5}{1.5. Extension of the approximately homogeneous piece.}\\
\hspace*{2em}\hyperlink{sec:1.6}{1.6. Properties of the extension.}\\
\hspace*{2em}\hyperlink{sec:1.7}{1.7. Formulation of the result. Theorem I.}\\
\hyperlink{chap:II}{Chapter II. Immediate consequences.}\\
\hspace*{2em}\hyperlink{sec:2.1}{2.1. Additivity and restricted weak continuity of $\TrM(A)$. Properties I, II.}\\
\hspace*{2em}\hyperlink{sec:2.2}{2.2. Definition of $\TrM(A)$ for all (not necessarily Hermitian) $A\in M$.}\\
\hspace*{4em}\hyperlink{sec:2.2}{Uniqueness. Properties III, IV.}\\
\hyperlink{chap:III}{Chapter III. The exact form of $\TrM(A)$ (for $\alpha\geq1$).}\\
\hspace*{2em}\hyperlink{sec:3.1}{3.1. The spectral form $\int_0^1\phi(\lambda)\,dE(\lambda)$ with $\DM(E(\lambda))=\lambda$.}\\
\hspace*{2em}\hyperlink{sec:3.2}{3.2. The maximum problems (A) and (B).}\\
\hspace*{2em}\hyperlink{sec:3.3}{3.3. Reducibility of a $g\in\cH$ by an $E\in M$.}\\
\hspace*{2em}\hyperlink{sec:3.4}{3.4. Decomposition of a $g$ from Theorem I.}\\
\hspace*{2em}\hyperlink{sec:3.5}{3.5. Construction of an exact $g$ for $\TrM(A)$. Removal of the condition}\\
\hspace*{4em}\hyperlink{sec:3.5}{$\alpha\geq1$. Unrestricted weak continuity of $\TrM(A)$. Theorems II, III, IV.}\\
\hyperlink{chap:IV}{Chapter IV. The isomorphism of $M$, $M'$, and $\cH$ (for $\alpha=1$).}\\
\hspace*{2em}\hyperlink{sec:4.1}{4.1. The isomorphism of $\cH$ and $\Q(M)$ with the help of a u.d.f.}\\
\hspace*{2em}\hyperlink{sec:4.2}{4.2. The correspondences of $\cH$, $\Q(M)$, and $\Q(M')$. Criteria for u.d.}\\
\hspace*{4em}\hyperlink{sec:4.2}{Equivalence of u.d. for $M$ and for $M'$. Isomorphism properties of}\\
\hspace*{4em}\hyperlink{sec:4.2}{these correspondences. Theorems V, VI.}\\
\hspace*{2em}\hyperlink{sec:4.3}{4.3. Analysis of $\Q(M)$, $\Q(M')$; their Hilbert-space character. Theorems}\\
\hspace*{4em}\hyperlink{sec:4.3}{VII, VIII, IX.}\\
\hspace*{2em}\hyperlink{sec:4.4}{4.4. Algebraic properties of $\Q(M)$. Roles of $M$, $M'$ in it. Properties}\\
\hspace*{4em}\hyperlink{sec:4.4}{$\mathrm I^0$, $\mathrm{II}^0$, $\mathrm{III}^0$, $\mathrm{IV}^0$, $\mathrm V^0$. Theorem X.}\\
\hspace*{2em}\hyperlink{sec:4.5}{4.5. Implication of a spatial isomorphism by an algebraic ring-isomorphism.}\\
\hspace*{4em}\hyperlink{sec:4.5}{Theorem XI.}\\
\hspace*{2em}\hyperlink{sec:4.6}{4.6. Description of further results.}\\
\hyperlink{chap:Appendix}{Appendix. The matrix-aspect of the $M$ of class $(\mathrm{II}_1)$.}

\chapterhead{I}{Approximate form of $\TrM(A)$ (for $\alpha\geq1$)}%
  {Chapter I. Approximate form of Tr M(A) (for α ≥ 1)}

\SectionStart{1.1}{1.1. The normalization}
1.1. We assume that we have given a factor $M$ in case $\mathrm{II}_1$. Now we normalize
$\DM$ and $\DMp$ so that $C=1$ (cf. \RO{}, pp.~179--182). This means (\RO{},
loc. cit.) that for every $f\in\cH$, $\DM(\mathfrak M_f^{M'})=\DMp(\mathfrak M_f^M)$ and that the range of $\DM$
is the interval $0\leq x\leq1$. Let the range of $\DMp$ be the interval $0\leq x\leq\alpha$, where
$0<\alpha\leq\infty$. In this part, we assume that $\alpha\geq1$.

Now $M'$ is a factor of class $\mathrm{II}_1$ or $\mathrm{II}_{\infty}$, by \RO{}, Theorem X. In the first case
the range of $\DMp$ is in the normalization (loc. cit.) the interval, $0\leq x\leq1$, and
$C=1/\alpha$ (since $\DMp=1/\alpha$ times its value in the previous paragraph). Since
$C\leq1$, we have in both cases by the discussion on page 182 of \RO{}, that
$\Delta_0=\Delta$. By \RO{}, Definition 10.1.1, this implies that there exists an $f$ such
that $\DM(\mathfrak M_f^{M'})=1$. Since we can pass from the normalization of this paragraph
to that of the previous paragraph by multiplying $\DMp$ by $\alpha$ and leaving
$\DM$ unchanged, we see that this holds even with our previous normalization.

Thus we have an $f$ such that $\DM(E_f^{M'})=1=\DM(1)$ or $\DM(1-E_f^{M'})=0$
which implies $1-E_f^{M'}=0$; that is, $E_f^{M'}=1$. We now assume that such an $f$
has been chosen and is held fixed for the following discussion.

\SectionStart{1.2}{1.2. The notions E ≥ λ, E ≤ λ, E ≥ₚ λ, E ≤ₚ λ}
1.2. With respect to this fixed $f$, we now define certain relations concerning
a projection $E\in M$.

\begin{statement}{definition:1.2.1}{Definition 1.2.1}
Let $E\ne0$ be a projection, $\in M$, $\theta$ a real number $\geq0$. The
relation $E>\theta$ ($E\geq\theta$, $E<\theta$, $E\leq\theta$) is said to hold if
$(Ef,f)>\theta\DM(E)$ ($(Ef,f)\geq\theta\DM(E)$, $(Ef,f)<\theta\DM(E)$,
$(Ef,f)\leq\theta\DM(E)$ respectively).

The relation $E\mathrel{\underset p\geq}\lambda$ ($E\mathrel{\underset p\leq}\lambda$) is said to hold for a projection $E\ne0$ if for every
$F\in M$ such that $F\leq E$, $F\ne0$, we have $F\geq\lambda$ ($F\leq\lambda$).
\end{statement}

\begin{statement}{lemma:1.2.1}{Lemma 1.2.1}
Let $\{E_i\}$ be a sequence of projections (finite or infinite) with
$E_i\in M$, $E_iE_j=\delta_j^iE_i$. Then if for every $i$, $E_i\leq\lambda$
($E_i\geq\lambda$), and if for some $i$, say $i_0$, $E_{i_0}<\lambda$
($E_{i_0}>\lambda$) we have $\sum_iE_i<\lambda$ ($\sum_iE_i>\lambda$).
\end{statement}

We have, for every $i$, $\lambda D(E_i)\leq(E_if,f)$, hence
$\lambda\sum_{i\ne i_0}D(E_i)\leq\sum_{i\ne i_0}(E_if,f)$ and
$\lambda D(E_{i_0})<(E_{i_0}f,f)$. These imply
$\lambda\sum_iD(E_i)<\sum_i(E_if,f)$ or
$\lambda D(\sum_iE_i)<(\sum_iE_if,f)$, which was to be shown.

\begin{statement}{lemma:1.2.2}{Lemma 1.2.2}
If $E\geq\lambda$ ($E\leq\lambda$), then there exists an $F$ such that
$E\geq F$, $F\mathrel{\underset p\geq}\lambda$ ($F\mathrel{\underset p\leq}\lambda$).
\end{statement}

By \StatementRef{definition:1.2.1}{Definition 1.2.1} we have $E\ne0$. Suppose there is no such $F$. Then there
must be an $E_1$, such that $E\geq E_1$, and $E_1<\lambda$. (Otherwise $E$ itself would be
such an $F$.) Now let $\Omega$ be the first ordinal number which belongs to a cardinal
number $\aleph>\aleph_0$. Let $\alpha$ be an ordinal number such that $\alpha<\Omega$. Let us suppose
that for all $\beta<\alpha$ an $E_\beta$ has been defined in such a way that $E_\beta\ne0$, $E_\beta\leq E$,
$E_\beta<\lambda$, $E_{\beta_1}E_{\beta_2}=\delta_{\beta_1,\beta_2}E_{\beta_1}$. Now if
$E-\sum_{\beta<\alpha}E_\beta=0$, let $E_{\alpha'}$ be undefined for $\alpha'\geq\alpha$. If, however,
$E-\sum_{\beta<\alpha}E_\beta\ne0$, by hypothesis, $E-\sum_{\beta<\alpha}E_\beta$ is not
$\mathrel{\underset p\geq}\lambda$, hence there exists an $E_0$, with $E_0\ne0$,
$E-\sum_{\beta<\alpha}E_\beta\geq E_0$, $E_0<\lambda$,
$E\geq E-\sum_{\beta<\alpha}E_\beta\geq E_0$. Thus if we let $E_\alpha=E_0$, we have, if
$E-\sum_{\beta<\alpha}E_\beta\ne0$, an $E_\alpha$ which is orthogonal to all previous $E_\beta$ and $E_\alpha<\lambda$.

Now $D(E)\geq\sum_{\beta<\alpha}D(E_\beta)$ for every $\alpha$. Since $D(E_\beta)>0$, this implies that
% Editorial note: the source's grammatically unusual singular construction is retained verbatim.
there is only denumerably many of the numbers $D(E_\beta)$. Hence for some
$\alpha<\Omega$, $E-\sum_{\beta<\alpha}E_\beta=0$. Inasmuch as $\alpha<\Omega$, we can re-index the $E_\beta$'s into a
finite or a simply infinite sequence. Now $E_1<\lambda$, $E_i\leq\lambda$, and since $E=\sum_iE_i$,
% Editorial note: source reads ``Lemma 1.1'' although the cited result is Lemma 1.2.1; wording retained verbatim and hyperlink targets the latter.
\StatementRef{lemma:1.2.1}{Lemma 1.1} implies that $E<\lambda$, a contradiction.

\begin{statement}{lemma:1.2.3}{Lemma 1.2.3}
If $E\mathrel{\underset p\geq}\lambda$ ($E\mathrel{\underset p\leq}\lambda$), then if $E\geq F$, $F\ne0$, $F\in M$, we have
$F\mathrel{\underset p\geq}\lambda$ ($F\mathrel{\underset p\leq}\lambda$).
\end{statement}

If $F_1$ is such that $F\geq F_1$, $F_1\ne0$ then $E\geq F\geq F_1$ and hence by
\StatementRef{definition:1.2.1}{Definition 1.2.1}, $F_1\geq\lambda$. This statement implies that $F\mathrel{\underset p\geq}\lambda$.

\SectionStart{1.3}{1.3. Selection of an approximately homogeneous piece}
1.3. Now $(1f,f)=\lVert f\rVert^2=\lVert f\rVert^2D(1)$. Now let
$\theta_0=\lVert f\rVert^2$ which is of course not zero for $\mathfrak M_f^M=\cH$ and this precludes
$f=0$. Thus 1 (the projection, not the number) is $\geq\theta_0$. By
\StatementRef{lemma:1.2.2}{Lemma 1.2.2}, there exists an $E\in M$, such that
$E\mathrel{\underset p\geq}\theta_0$. Now let $\lambda$ be the least upper bound of the numbers $\theta$ such that
$E\mathrel{\underset p\geq}\theta$. $\lambda$ must be less than $(Ef,f)/\DM(E)$ and since $\DM(E)>0$, it must be finite.
Let $\epsilon$ be any number $>0$. $E$ is not $\mathrel{\underset p\geq}\lambda+\epsilon$. Hence there is an $E_1$, such that
$E\geq E_1$, and $E_1\leq\lambda+\epsilon$. \StatementRef{lemma:1.2.2}{Lemma 1.2.2} now implies that there is an $E_2$ with
$E_1\geq E_2$ and $E_2\mathrel{\underset p\leq}\lambda+\epsilon$.

Now if $F$ is such that $E\geq F$, $F\in M$, $F\ne0$, we have $F\geq\theta$ for all $\theta<\lambda$, which
means of course that $F\geq\lambda$ and $E\mathrel{\underset p\geq}\lambda$. Now since $E\geq E_1\geq E_2$,
\StatementRef{lemma:1.2.3}{Lemma 1.2.3} implies that $E_2\mathrel{\underset p\geq}\lambda$. Thus for some fixed $\lambda>0$, given any $\epsilon>0$, we can find a
non-zero $E_2\in M$, such that for all $F\in M$ with the property $E_2\geq F$, we have
\[
  (\lambda+\epsilon)\DM(F)\geq(Ff,f)\geq\lambda\DM(F).
\]

% Editorial note: source reads $\lambda^{1/2}f$; retained verbatim despite the apparent scaling misprint.
Now if we let $f$ be $\lambda^{1/2}f$ above, $\epsilon=\lambda(K-1)$, $E_2=E$, we see that

\begin{statement}{lemma:1.3.1}{Lemma 1.3.1}
To every $K>1$ there exists an $f$ and $E\in M$, $E\ne0$, such that
for every $F\in M$ with the property $E\geq F$,
\[
  K\DM(F)\geq(Ff,f)\geq\DM(F).
\]
\end{statement}

\SectionStart{1.4}{1.4. Properties of such a piece}
1.4. We can now show

\begin{statement}{lemma:1.4.1}{Lemma 1.4.1}
Let $K$, $f$, and $E$ be as in \StatementRef{lemma:1.3.1}{Lemma 1.3.1}. Let $A\in M$ be a positive
definite self-adjoint operator such that $EA=AE=A$. Then
\[
  K\TrM(A)\geq(Af,f)\geq\TrM(A).
\]
\end{statement}

Let $c$ be the bound of $A$, $E(\lambda)$ the resolution of the identity corresponding
to $A$. Now by \RO{}, pp.~212--213,
\[
  \TrM(A)=\int_0^c\lambda\,dD(E(\lambda));
  \qquad
  (Af,f)=\int_0^c\lambda\,d(E(\lambda)f,f).
\]

Now since $AE=A$, $E(0)\geq1-E$, and since $A$ commutes with $E$, $E$ commutes
with all $E(\lambda)$ (cf. \ROcite{15}, Theorem~8.2). Now
$E(\lambda)=E(\lambda)E+E(\lambda)(1-E)$. Since $E(\lambda)$ and $E$ commute,
$E(\lambda)E=F(\lambda)$ is a projection and similarly $E(\lambda)(1-E)$ is a projection.
For $\lambda\geq0$, we have $E(\lambda)(1-E)\leq1-E$ and also
$E(\lambda)(1-E)\geq E(0)(1-E)\geq(1-E)^2=1-E$ and thus
$E(\lambda)(1-E)=1-E$. So for $\lambda\geq0$, $E(\lambda)=F(\lambda)+(1-E)$.
Now $F(\lambda)(1-E)=E(\lambda)E(1-E)=0$, hence by \RO{}, Definition 8.2.1,
$D(E(\lambda))=D(F(\lambda))+D(1-E)$.
Thus we have
\begin{align*}
  \TrM(A)
   &=\int_0^c\lambda\,dD(E(\lambda))
    =\int_0^c\lambda\,dD(F(\lambda))
      +\int_0^c\lambda\,dD(1-E) \\
   &=\int_0^c\lambda\,dD(F(\lambda))
\end{align*}
and
\begin{align*}
  (Af,f)
   &=\int_0^c\lambda\,d(E(\lambda)f,f) \\
   &=\int_0^c\lambda\,d(F(\lambda)f,f)
     +\int_0^c\lambda\,d((1-E)f,f) \\
   &=\int_0^c\lambda\,d(F(\lambda)f,f).
\end{align*}

But \StatementRef{lemma:1.3.1}{Lemma 1.3.1} now implies that if $0\leq\alpha<\beta\leq c$, then
\[
  K\DM(F(\beta)-F(\alpha))\geq((F(\beta)-F(\alpha))f,f)\geq\DM(F(\beta)-F(\alpha))
\]
which by the definition of the Riemann-Stieltjes integral yields that
\[
  K\TrM(A)\geq(Af,f)\geq\TrM(A).
\]

\begin{statement}{lemma:1.4.2}{Lemma 1.4.2}
For a projection $E\ne0$, $E\mathrel{\underset p\geq}\theta$
($E\mathrel{\underset p\leq}\theta$) is equivalent to the statement that for all positive definite $A\in M$
such that $EA=AE=A$,
\[
  (Af,f)\geq\theta\TrM(A)\quad(\leq\theta\TrM(A)).
\]
\end{statement}

Since $\TrM(F)=\DM(F)$, $F$ a projection, the last statement implies the
first. The converse is shown by a proof similar to that of
\StatementRef{lemma:1.4.1}{Lemma 1.4.1}.

\begin{statement}{lemma:1.4.3}{Lemma 1.4.3}
Let $A\in M$ be such that $EA=AE=A$, then
$\lVert A^*f\rVert^2\leq K\lVert Af\rVert^2$, if $E$, $f$, $K$ are as in
\StatementRef{lemma:1.3.1}{Lemma 1.3.1}.
\end{statement}

$A^*A$ is positive definite and $\in M$. Furthermore $AE=EA=A$ implies
$EA^*=A^*E=A^*$, and thus $EA^*A=A^*A=A^*AE$. Hence
\StatementRef{lemma:1.4.1}{Lemma 1.4.1} applies to $A^*A$ and we have
\EquationTarget{alpha}
\begin{equation*}
  K\TrM(A^*A)\geq(A^*Af,f)\geq\TrM(A^*A). \tag{$\alpha$}
\end{equation*}
Using the canonical decomposition (\RO{}, Definition 4.4.1) for $A^*$ we have
$A^*=UB$, where $U$ may be taken as unitary. (In the finite cases, we see
(cf. \RO{}, Lemma 16.1.1) that there exists a partially isometric $V$ with initial
set $(f;A^*f=0)$ and final set $(f;Af=0)$. Now if $W$ is as in \RO{}, Definition
4.4.1, for $A=A^*$, let $U=W+V$.) $B$ is self-adjoint and equals $(AA^*)^{1/2}$. Now
$A=BU^*=BU^{-1}$ and hence $A^*A=UBBU^{-1}=UB^2U^{-1}=UAA^*U^{-1}$. Hence
\EquationTarget{beta}
\begin{equation*}
  \TrM(A^*A)=\TrM(UAA^*U^{-1})=\TrM(AA^*). \tag{$\beta$}
\end{equation*}

Substituting $A^*$ for $A$ in our previous result we have
\EquationTarget{gamma}
\begin{equation*}
  K\TrM(AA^*)\geq(AA^*f,f)\geq\TrM(AA^*). \tag{$\gamma$}
\end{equation*}
% Editorial note: source cites $(\phi)$ here although the preceding displayed formula is labelled $(\gamma)$; retained verbatim and hyperlink targets $(\gamma)$.
Combining \EquationRef{alpha}{$(\alpha)$}, \EquationRef{beta}{$(\beta)$}, and
\EquationRef{gamma}{$(\phi)$}, we obtain
\[
  K(A^*Af,f)\geq K\TrM(A^*A)=K\TrM(AA^*)\geq(AA^*f,f).
\]
Since $(A^*Af,f)=(Af,Af)=\lVert Af\rVert^2$,
$(AA^*f,f)=(A^*f,A^*f)=\lVert A^*f\rVert^2$; this is the desired inequality.

\SectionStart{1.5}{1.5. Extension of the approximately homogeneous piece}
1.5. Now if $E$ is as in \StatementRef{lemma:1.3.1}{Lemma 1.3.1}, let $n$ be the smallest integer such
that $1/n\leq\DM(E)$. Then if $E^0\leq E$ is such that $\DM(E^0)=1/n$,
\StatementRef{lemma:1.3.1}{Lemma 1.3.1} will hold with $E^0$ in place of $E$. Now $f$ was chosen in such a manner that
$\mathfrak M_f^{M'}=\cH$. This implies that $\mathfrak M_{E^0f}^{M'}$ is the range of $E^0$, because since the set
$(A'f;A'\in M')$ is dense in $\cH$, the set $(E^0A'f;A'\in M')=(A'E^0f;A'\in M')$ must be
dense in the range of $E^0$, or $E^0=E_{E^0f}^{M'}$. Since $C=1$,
$\DMp(E_{E^0f}^{M})=\DM(E_{E^0f}^{M'})=\DM(E^0)=1/n$.

Now let $E^0=E_1$, $E_{E^0f}^{M}=E_1'$. Since $\DM(E_1)=1/n$, $\DM(1)=1$, we can find $n$
projections $\{E_j\}$, $j=1,\cdots,n$ with the first equal to $E_1$ such that $E_j\in M$,
$\sum_{j=1}^nE_j=1$, $E_iE_j=0$ if $i\ne j$, $\DM(E_j)=1/n$. Since
$\DMp(E_1')=1/n$, $\DMp(1)=\alpha\geq1$, there exist $n$ projections $E_j'\in M'$,
$j=1,\cdots,n$, with the first equal to $E_1'$ and such that $E_i'E_j'=0$ if $i\ne j$ and
$\DMp(E_j')=1/n$. Since $\DM(E_1)=\DM(E_j)$, there is a partially isometric
operator $W_j\in M$ with initial set the range of $E_1$ and final set the range of
$E_j$ (cf. \RO{}, Definition 4.3.1). Similarly we have a $W_j'\in M'$ with initial
set the range of $E_1'$ and final set $E_j'$. The relations
$W_j^*W_j=E_1$, $W_jW_j^*=E_j$, $W_j'^*W_j'=E_1'$ and
$W_j'W_j'^*=E_j'$ hold, also $W_jE_1=W_j$, $E_jW_j=W_j$,
$E_1W_j^*=W_j^*$, $W_j^*E_j=W_j^*$, $W_j'E_1'=W_j'$,
$E_j'W_j'=W_j'$, $E_1'W_j'^*=W_j'^*$, $W_j'^*E_j'=W_j'^*$
(cf. \RO{}, Lemma 4.3.1).

Let $W_j'W_jE_1f=f_j$, $E_1f=f_1$. Then
$E_jf_j=E_jW_j'W_jE_1f=W_j'E_jW_jE_1f=W_j'W_jE_1f=f_j$ and
$E_j'f_j=E_j'W_j'W_jE_1f=W_j'W_jE_1f=f_j$. Thus
$E_jf_j=f_j=E_j'f_j$.

Let $g=\sum_{j=1}^nf_j$, then
$E_ig=E_i\sum_{j=1}^nf_j=E_i\sum_{j=1}^nE_jf_j
=\sum_{j=1}^nE_iE_jf_j=E_if_i=f_i$.
Similarly $E_i'g=f_i$. Now for any $A\in M$ if either $i\ne j$ or $k\ne l$, then
$E_iAE_kg$ is orthogonal to $E_jAE_lg$. For inasmuch as $E_kg=f_k=E_k'g$, we have,
\begin{align*}
 (E_iAE_kg,E_jAE_lg)
  &=(E_jE_iAE_kg,AE_lg)
    =\delta_j^i(E_iAE_kg,AE_lg)\\
  &=\delta_j^i(E_iAE_k'g,AE_l'g)
    =\delta_j^i(E_k'E_iAg,E_l'Ag)\\
  &=\delta_j^i(E_l'E_k'E_iAg,Ag)
    =\delta_j^i\delta_l^k(E_k'E_iAg,Ag).
\end{align*}

These results imply that if $A\in M$, then
\begin{align*}
 \lVert Ag\rVert^2
 &=\left\lVert\left(\sum_{i=1}^nE_i\right)
        A\left(\sum_{j=1}^nE_j\right)g\right\rVert^2\\
 &=\left\lVert\sum_{i=1}^n\sum_{j=1}^nE_iAE_jg\right\rVert^2
  =\sum_{i=1}^n\sum_{j=1}^n\lVert E_iAE_jg\rVert^2\\
 &=\sum_{i=1}^n\sum_{j=1}^n\lVert E_iAE_jE_jg\rVert^2
  =\sum_{i=1}^n\sum_{j=1}^n\lVert E_iAE_jf_j\rVert^2\\
 &=\sum_{i=1}^n\sum_{j=1}^n\lVert E_iAE_jW_j'W_jf_1\rVert^2\\
 &=\sum_{i=1}^n\sum_{j=1}^n\lVert W_j'E_iAE_jW_jf_1\rVert^2\\
 &=\sum_{i=1}^n\sum_{j=1}^n\lVert W_j'E_1'E_iAE_jW_jf_1\rVert^2\\
 &=\sum_{i=1}^n\sum_{j=1}^n\lVert E_1'E_iAE_jW_jf_1\rVert^2\\
 &=\sum_{i=1}^n\sum_{j=1}^n\lVert E_iAE_jW_jE_1'f_1\rVert^2\\
 &=\sum_{i=1}^n\sum_{j=1}^n\lVert E_iAE_jW_jf_1\rVert^2\\
 &=\sum_{i=1}^n\sum_{j=1}^n\lVert W_i^*E_iAE_jW_jf_1\rVert^2\\
 &=\sum_{i=1}^n\sum_{j=1}^n\lVert W_i^*E_iAE_jW_jE_1f\rVert^2\\
 &=\sum_{i=1}^n\sum_{j=1}^n\lVert W_i^*E_iAE_jW_jf\rVert^2,
\end{align*}
remembering that $W_j'$ is isometric on the range of $E_1'$, $W_i^*$ on the range of $E_i$
and $W_jE_1=W_j$. Substituting $A^*$ for $A$ and interchanging $i$ and $j$ we also have
\begin{align*}
 \lVert A^*g\rVert^2
 &=\sum_{i=1}^n\sum_{j=1}^n\lVert W_j^*E_jA^*E_iW_if\rVert^2\\
 &=\sum_{i=1}^n\sum_{j=1}^n\lVert(W_i^*E_iAE_jW_j)^*f\rVert^2.
\end{align*}

But recalling the properties of $W_j$ and $W_i^*$ we have
\[
 E_1W_i^*E_iAE_jW_j=W_i^*E_iAE_jW_j=W_i^*E_iAE_jW_jE_1.
\]
Thus \StatementRef{lemma:1.4.3}{Lemma 1.4.3} with $A=W_i^*E_iAE_jW_j$ yields
$\lVert(W_i^*E_iAE_jW_j)^*f\rVert^2
\allowbreak\leq K\lVert W_i^*E_iAE_jW_jf\rVert^2$
and with this, we obtain from the above equations for $\lVert Ag\rVert^2$ and
$\lVert A^*g\rVert^2$ that $\lVert A^*g\rVert^2\leq K\lVert Ag\rVert^2$.

This result may be stated as follows.

\begin{statement}{lemma:1.5.1}{Lemma 1.5.1}
If $K$ is any number $>1$, there exists a $g\in\cH$, $g\ne0$ such that for
every $A\in M$, $\lVert A^*g\rVert^2\leq K\lVert Ag\rVert^2$. Obviously we may assume that
$\lVert g\rVert=1$.
\end{statement}

\SectionStart{1.6}{1.6. Properties of the extension}
1.6. We now have

\begin{statement}{lemma:1.6.1}{Lemma 1.6.1}
Let $K$ and $g$ be as in \StatementRef{lemma:1.5.1}{Lemma 1.5.1}. Then if $E$ and $F$, $\in M$, are
two projections such that $\DM(E)=\DM(F)$, then
$K^{-1}(Fg,g)\leq(Eg,g)\leq K(Fg,g)$.
\end{statement}

Inasmuch as $\DM(E)=\DM(F)$, there exists a partially isometric operator
$W$ such that $W^*W=F$, $WW^*=E$, $W\in M$ (cf. \RO{}, Definition 8.2.1, Definition
6.1.1, and Lemma 4.3.1). Now
$(Fg,g)=(W^*Wg,g)=(Wg,Wg)=\lVert Wg\rVert^2$,
$(Eg,g)=(WW^*g,g)=(W^*g,W^*g)=\lVert W^*g\rVert^2$.
\StatementRef{lemma:1.5.1}{Lemma 1.5.1} with $A=W$ implies $(Eg,g)\leq K(Fg,g)$.
With $A=W^*$, the same lemma implies $(Fg,g)\leq K(Eg,g)$ or
$K^{-1}(Fg,g)\leq(Eg,g)$. We have now shown our lemma.

Suppose $E\in M$ is such that $\DM(E)=1/m$, where $m$ is an integer. Then
there exist $m-1$ projections $E_2,\cdots,E_m$, such that when we let $E_1=E$, we
have $\sum_{j=1}^mE_j=1$, $E_iE_j=0$, for $i\ne j$, $E_j\in M$, $\DM(E_j)=1/m$.
Now returning to the $g$ of \StatementRef{lemma:1.5.1}{Lemmas 1.5.1} and
\StatementRef{lemma:1.6.1}{1.6.1}, we have
\[
 1=\lVert g\rVert^2=(g,g)=\left(\sum_{j=1}^mE_jg,g\right)
   =\sum_{j=1}^m(E_jg,g).
\]
Let $(E_1g,g)=\alpha$. Then \StatementRef{lemma:1.6.1}{Lemma 1.6.1} implies that
$K\alpha\geq(E_jg,g)\geq K^{-1}\alpha$, for every $j$. Summing over $j$, gives
$Km\alpha\geq1\geq K^{-1}m\alpha$ or $K\alpha\geq1/m\geq K^{-1}\alpha$ which is the same as
\EquationTarget{plus}
\begin{equation*}
 K(Eg,g)\geq\DM(E)\geq K^{-1}(Eg,g), \tag{$+$}
\end{equation*}
since $\DM(E)=1/m$.

Let us study the class $\Sigma$ of $E$'s in $M$ which satisfy the equation
\EquationRef{plus}{$(+)$}. Now if $F_1,\cdots,F_q$ or $F_1,F_2,\cdots$ satisfy
\EquationRef{plus}{$(+)$} and $F_iF_j=0$, $i\ne j$, then
$\sum_{j=1}^qF_j$ or $\sum_{j=1}^{\infty}F_j$ also satisfies
\EquationRef{plus}{$(+)$}. But if $F\in M$ is such that $\DM(F)=q/p$,
$F=\sum_{j=1}^qF_j$ for mutually orthogonal $F_j$'s $\in M$ with
$\DM(F_j)=1/p$ and hence satisfying \EquationRef{plus}{$(+)$}. Thus if
$\DM(F)=q/p$, $F$ satisfies \EquationRef{plus}{$(+)$}.

Let $E$ be any projection of $M$, $\DM(E)=\alpha$. Let $\{\alpha_i\}$, $\alpha_i\geq0$ be a sequence
of rational numbers, $\sum_{i=1}^{\infty}\alpha_i=\alpha$. Then there exists a set of mutually orthogonal
projections $\{E_i\}$ such that $E_i\leq E$, $E_i\in M$, $\DM(E_i)=\alpha_i$. To show this,
suppose $E_1,\cdots,E_{i-1}$ have been chosen with these properties. Then
$\DM\left(E-\sum_{j=1}^{i-1}E_j\right)=\sum_{j=i}^{\infty}\alpha_j\geq\alpha_i$.
Hence an $E_i$ can be chosen in such a way that $E_i\in M$ is
$\leq E-\sum_{j=1}^{i-1}E_j$ with $\DM(E_i)=\alpha_i$. Now
$\DM\left(E-\sum_{i=1}^{\infty}E_i\right)
=\alpha-\sum_{i=1}^{\infty}\alpha_i=0$
or $E=\sum_{i=1}^{\infty}E_i$. Since each $E_i$ satisfies
\EquationRef{plus}{$(+)$}, $E$ does too. Since $E$ is arbitrary we have

\begin{statement}{lemma:1.6.2}{Lemma 1.6.2}
Let $K$ and $g$ be as in \StatementRef{lemma:1.5.1}{Lemma 1.5.1}, then for every $E\in M$
\[
 K(Eg,g)\geq\DM(E)\geq K^{-1}(Eg,g).
\]
\end{statement}

\SectionStart{1.7}{1.7. Formulation of the result. Theorem I}
1.7. From \StatementRef{lemma:1.6.2}{Lemma 1.6.2} we can conclude, using
\StatementRef{lemma:1.4.2}{Lemma 1.4.2}, that if $A\in M$ is positive definite, then
\[
 K(Ag,g)\geq\TrM(A)\geq K^{-1}(Ag,g).
\]

We can state our result as follows.

\begin{statement}{theorem:I}{Theorem I}
Let $M$ be a factor in case $\mathrm{II}_1$. Let $M$ and $M'$ be such that
when $\DM$ and $\DMp$ are normalized in such a way that the range of $\DM$ is the
interval $(0,1)$ and $C=1$ (cf. \SectionRef{1.1}), then the range of $\DMp$ is an interval $(0,\alpha)$
% Editorial note: the source's consecutive ``then ... Then'' construction is retained verbatim.
with $0<\alpha\leq\infty$, then $\alpha\geq1$. Then to every $K>1$, there exists a $g\in\cH$ such that for
every positive definite $A\in M$,
\[
 K(Ag,g)\geq\TrM(A)\geq K^{-1}(Ag,g).
\]
\end{statement}

\chapterhead{II}{Immediate consequences}{Chapter II. Immediate consequences}

\SectionStart{2.1}{2.1. Additivity and restricted weak continuity of Tr M(A). Properties I, II}
2.1. We are now able to show that $\TrM$ has the following two properties
when $M$ is in a finite case (cf. \RO{}, Theorem VIII); that is, when $M$ is in a
case $\mathrm I_n$ ($n=1,2,\cdots$) or $\mathrm{II}_1$.

\begin{statement}{property:I}{Property I}
For all Hermitian $A$ and $B\in M$
\[
 \TrM(A+B)=\TrM(A)+\TrM(B).
\]
\end{statement}

\begin{statement}{property:II}{Property II}
$\TrM(A)$ is weakly continuous if $A$ is subject to either of these conditions:
\begin{enumerate}[label=(\roman*)]
 \item $A$ is uniformly bounded (that is, $\opnorm{A}\leq D$ for some fixed $D$);
 \item $A$ is definite.
\end{enumerate}
\end{statement}

These properties are obviously independent of the normalization of $\DM(E)$
and $\TrM(A)$.

If $M$ is in cases $\mathrm I_n$, $n=1,2,\cdots$, these are of course well known results
about $\TrM(A)$ (cf. \RO{}, p.~220). So we may assume $M$ to be in case $\mathrm{II}_1$.

Then we have in the normalization of \RO{}, Theorem VIII, three possibilities
for the factorization $M$, $M'$: $\mathrm{II}_1$, $\mathrm{II}_{\infty}$; $\mathrm{II}_1$, $\mathrm{II}_1$ with $C\leq1$;
$\mathrm{II}_1$, $\mathrm{II}_1$ with $C\geq1$. The two first ones correspond to $\alpha\geq1$ in the normalization of
\SectionRef{1.1}, and since we shall only use \StatementRef{theorem:I}{Theorem I}, from \SectionRef{1.1}, we may treat these two
cases together.

We therefore consider first the conditions under which \StatementRef{theorem:I}{Theorem I} holds:
the normalization of \SectionRef{1.1}, and $\alpha\geq1$. We obtain an equivalent statement of
\StatementRef{theorem:I}{Theorem I} with $A\in M$, Hermitian and not necessarily definite.

Let $K$, $g$ be as in \StatementRef{theorem:I}{Theorem I}, then
$K^{-1}(g,g)=K^{-1}\lVert g\rVert^2\leq\TrM(1)=1$ or
$\lVert g\rVert^2\leq K$. Furthermore if $A\in M$ is Hermitian, then $A=A_1-A_2$, where $A_1$ and $A_2$
are both positive definite and $A_1A_2=0=A_2A_1$. (Take
$A_1=\tfrac12(A+\lvert A\rvert)$, $A_2=\tfrac12(\lvert A\rvert-A)$
(cf. \ROcite{15}, Chapter VI, or \ROcite{19}, p.~203 ff.).) Since $A_1$ and $A_2$
commute, $\TrM(A)=\TrM(A_1)-\TrM(A_2)$ by \RO{}, Lemma 15.3.4. It is a consequence
% Editorial note: source prints the Arabic numeral in Theorem 1 here; retained verbatim.
of \StatementRef{theorem:I}{Theorem 1} that, if $\epsilon=K-1$,
\[
 \lvert\TrM(A_1)-(A_1g,g)\rvert\leq\epsilon(A_1g,g);
 \qquad
 \lvert\TrM(A_2)-(A_2g,g)\rvert\leq\epsilon(A_2g,g).
\]
These with the above equation for $\TrM(A)$ imply
\[
 \lvert\TrM(A)-(Ag,g)\rvert
 \leq\epsilon\bigl((A_1g,g)+(A_2g,g)\bigr).
\]
But the bounds of $A_1$ and $A_2$ are each not more than that of $A$. Also
$\lVert g\rVert^2\leq1+\epsilon$. So if $A$ has the bound $D$, we have
% Editorial note: source refers to this as (1) below but prints no visible tag here;
% the target is intentionally hidden while preserving a working link.
\EquationTarget{1}
\[
 \lvert\TrM(A)-(Ag,g)\rvert\leq2\epsilon(1+\epsilon)D.
\]

Let $A$, $B$, $A+B$ have the respective bounds $D_1$, $D_2$, $D_3$. Substituting in
\EquationRef{1}{(1)} each of these operators, inasmuch as
$((A+B)g,g)=(Ag,g)+(Bg,g)$, we get
\[
 \lvert\TrM(A+B)-\TrM(A)-\TrM(B)\rvert
 \leq2\epsilon(1+\epsilon)(D_1+D_2+D_3),
\]
which since $\epsilon$ may be taken arbitrarily small, implies
\StatementRef{property:I}{Property I}.

We turn our attention to \StatementRef{property:II}{Property II}, (i). Consider the set $\Sigma$ of all
Hermitian $A\in M$, with a bound less than or equal to a fixed $D$. We will
show that to every $A\in\Sigma$ and to every $\eta>0$, there exists a weak neighborhood
of $A$, $U(A;g;g;\eta/3)$ such that for all $B\in\Sigma$ and
$B\in U(A;g;g;\eta/3)$ (i.e., $\lvert((B-A)g,g)\rvert<\eta/3$) we have
$\lvert\TrM(B)-\TrM(A)\rvert<\eta$. Letting $\epsilon$ be such that
$\epsilon(1+\epsilon)D\leq\eta/3$, as above we can find a $g$ such that for $A$ and every $B\in\Sigma$,
$\lvert\TrM(A)-(Ag,g)\rvert\leq\epsilon(1+\epsilon)D\leq\eta/3$,
$\lvert\TrM(B)-(Bg,g)\rvert\leq\eta/3$.
In particular for $B\in U(A;g;g;\eta/3)$ (cf. above) which means that we also have
$\lvert((B-A)g,g)\rvert<\eta/3$, these inequalities imply that
$\lvert\TrM(A)-\TrM(B)\rvert<\eta$.

Thus we have shown the weak neighborhood continuity of $\TrM(A)$ with
of course a restriction. From this the sequential continuity follows immediately.
(If $A_n\to A$, then the $A_n$ are uniformly bounded and for any $\eta>0$,
almost all the $A_n$ must be in the above given neighborhood.) In this situation
the two kinds of continuity are equivalent (cf. \ROcite{18}, pp.~383--384), but we
will use only the sequential.

$\TrM(A)$ is also weakly continuous when considered only for the positive
definite $A\in M$. For suppose $A$ is such and an $\epsilon>0$ is given. Let the $K$ of
\StatementRef{theorem:I}{Theorem I} be chosen in such a way that $1<K\leq2$, and
$(K-1)\TrM(A)<\epsilon/5$, and let $g$ correspond to this $K$ as there. Let $\eta$ be chosen in such a way
that $\eta\leq\epsilon/5$. Then if a positive definite $B\in M$ is in $U(A;g;g;\eta)$, we have
$\lvert\TrM(A)-\TrM(B)\rvert<\epsilon$ (because
$\lvert((B-A)g,g)\rvert<\eta\leq\epsilon/5$,
$\lvert\TrM(A)-(Ag,g)\rvert\leq(K-1)\TrM(A)<\epsilon/5$, and, in addition,
$\lvert\TrM(B)-(Bg,g)\rvert\leq(K-1)(Bg,g)
\leq(K-1)((Ag,g)+\eta)
\leq(K-1)(\TrM(A)+\epsilon/5+\eta)\leq3\epsilon/5$).

Thus we have settled the factorizations $\mathrm{II}_1$, $\mathrm{II}_{\infty}$ and
$\mathrm{II}_1$, $\mathrm{II}_1$ with $C\leq1$. But we still must consider the case $\mathrm{II}_1$, $\mathrm{II}_1$ with $C>1$.
% Editorial note: the source joins these clauses with a comma and omits a finite verb before $E_m$; retained verbatim.
Let $m$ be an integer such that $C/m\leq1$, $E_m$ an $m$-dimensional unitary space.
Form $E_m\otimes\cH$ (cf. \RO{}, Theorem I). Let $N$ be the ring of operators on $E_m$,
$N'$ the corresponding set of operators on $E_m\otimes\cH$ (cf. \RO{}, Lemma 2.3.1 and
Lemma 2.3.2), $B$ the set of all operators in $\cH$, $B^{(2)}$ the corresponding set in
$E_m\otimes\cH$.

Under the above correspondence of $B$ and $B^{(2)}$, $M\sim M^{(2)}$,
$M'\sim M'^{(2)}$. The first and hence each of the others is a full ring isomorphism
(\RO{}, Lemma 2.3.4 and \ROcite{22}, \S4). Thus $M^{(2)}$ and $M'^{(2)}$ are in case
$\mathrm{II}_1$ (\RO{}, Theorem IX). Now let $\phi_1,\cdots,\phi_m$ be a complete orthonormal
set in $E_m$ and let us think of the set up of \RO{}, \S2.4. By \RO{}, Lemma 2.4.5, we
have $M^{(2)\prime}=R(N^{(1)},M'^{(2)})$ and by \RO{}, Lemma 11.5.2,
$R(N^{(1)},M'^{(2)})$ is in case $\mathrm{II}_1$ since as we have mentioned before,
$M'^{(2)}$ is in this case.

Thus the coupled factorization, $M^{(2)}$, $M^{(2)\prime}$, is in case $\mathrm{II}_1$, $\mathrm{II}_1$ and
$M^{(2)}$ is fully ring isomorphic to $M$. Furthermore such an isomorphism preserves
$\DM$ in the standard normalization (\RO{}, Theorem IX) and hence the trace. Now
suppose we have shown that for $M^{(2)}$, $M^{(2)\prime}$ we have $C^{(2)}\leq1$. Then
$M^{(2)}$ has \StatementRef{property:I}{Properties I} and
% Editorial note: the source's grammatically irregular wording is retained verbatim.
\StatementRef{property:II}{II}, by the above and $M$ must have them to be the
isomorphism. (As $m$ is finite, even weak operator topology is conserved under
the mapping $B\sim B^{(2)}$.)

Therefore it remains to show that $C^{(2)}\leq1$. Consider the closed linear manifold,
$(\phi_1\otimes f;f\in\cH)=\text{(say) }\mathfrak M$, in $E_m\otimes\cH$. The correspondence of
$\mathfrak M$ and $\cH$ given by $\phi_1\otimes f\sim f$ is unitary. Under it,
$M_{\mathfrak M}^{(2)}$ and $R(N^{(1)},M'^{(2)})_{\mathfrak M}$ (cf. \RO{}, Definition 11.3.1)
are unitarily equivalent to $M$ and $M'$ (cf. \RO{}, \S2.4). Thus the $C$ for $M$, $M'$
is the same as that of $M_{\mathfrak M}^{(2)}$ and $R(N^{(1)},M'^{(2)})_{\mathfrak M}$.
Then from the proof of \RORef{Lemma 11.4.3}, we can conclude that
$C=(D_{M^{(2)}}(\cH)/D_{M^{(2)\prime}}(\mathfrak M))C^{(2)}
=mC^{(2)}$
or $C^{(2)}=C/m\leq1$.

\SectionStart{2.2}{2.2. Definition of Tr M(A) for all (not necessarily Hermitian) A in M. Uniqueness. Properties III, IV}
2.2. Previously $\TrM(A)$ has only been considered for Hermitian $A\in M$. We
now define it for all $A\in M$.

\begin{statement}{definition:2.2.1}{Definition 2.2.1}
If $A\in M$, then $A=B+iC$, where
$B=\tfrac12(A+A^*)$, $C=-\tfrac{1}{2}i(A-A^*)$ are Hermitian, and we define
$\TrM(A)=\TrM(B)+i\TrM(C)$.
\end{statement}

$\TrM(A)$ has the following property.

\begin{statement}{property:III}{Property III}
\EquationTarget{property:III:star}$(*)$ $\TrM(A)$ agrees with the previous definition of $\TrM(A)$ if $A$ is Hermitian.

\EquationTarget{property:III:starstar}$(*\,*)$ For uniformly bounded or for definite $A$'s, $\TrM(A)$ is weakly continuous.
\begin{enumerate}[label=(\roman*)]
 \item $\TrM(1)=1$;
 \item $\TrM(\rho A)=\rho\TrM(A)$, $\rho$ complex;
 \item $\TrM(A+B)=\TrM(A)+\TrM(B)$;
 \item $\TrM(A)\geq0$, if $A$ is definite;
 \item[(iv)$'$] If $A$ is definite and $\ne0$, $\TrM(A)>0$;
 \item $\TrM(A^*)=\overline{\TrM(A)}$;
 \item $\TrM(AB)=\TrM(BA)$;
 \item[(vi)$'$] $\TrM(U^{-1}AU)=\TrM(A)$, if $U$ is unitary.
\end{enumerate}
\end{statement}

Now, \EquationRef{property:III:star}{$(*)$} and (v) are immediate consequences of the definition. (i), (iv),
and (iv)$'$ follow from \EquationRef{property:III:star}{$(*)$} and \RO{}, Theorem XIII. $A^*$ is a weakly continuous
additive function of $A$. Thus $B$ and $C$ are too, and we may infer (iii) and
\EquationRef{property:III:starstar}{$(*\,*)$} from \StatementRef{property:I}{Properties I} and
\StatementRef{property:II}{II} respectively. To show (ii) we first note that
if $A$ is Hermitian by \StatementRef{definition:2.2.1}{Definition 2.2.1},
$\TrM(iA)=i\TrM(A)$ and if $a$ is real by \EquationRef{property:III:star}{$(*)$} $\TrM(aA)=a\TrM(A)$. Then (ii)
can be shown by a calculation involving (iii).

To show (vi) suppose $A$ and $B\in M$ are Hermitian. Let $C=A+iB$,
$C^*=A-iB$. Then $\TrM(CC^*)=\TrM(C^*C)$ (by \EquationRef{property:III:star}{$(*)$}, cf. proof of
\StatementRef{lemma:1.4.3}{Lemma 1.4.3}, equation \EquationRef{beta}{$(\beta)$}).
Substituting for $C$, expanding the products according to the distributive law
and collecting terms, we get
$2i(\TrM(AB)-\TrM(BA))=0$ or $\TrM(AB)=\TrM(BA)$. From this we can show
(vi) in general by using \StatementRef{definition:2.2.1}{Definition 2.2.1},
(ii) and (iii). (vi)$'$ follows from (vi) by letting $A=U^{-1}A$, $B=U$ and
using the associative law.

\begin{statement}{property:IV}{Property IV}
(i)--(iv), (v), (vi)$'$ of \StatementRef{property:III}{Property III} characterize
$\TrM(A)$ uniquely.
\end{statement}

Let $\Tr'(A)$ have the listed properties. Then if $A$ is Hermitian, $\Tr'(A)$
is real by (v). This and (i)--(iv), (vi)$'$ imply that for $A\in M$ and Hermitian
$\Tr'(A)=\TrM(A)$ (\RO{}, Theorem XIII). By (ii) and (iii), we then have that
for an arbitrary $C\in M$, $C=A+iB$, $A$ and $B$ Hermitian,
$\Tr'(C)=\Tr'(A+iB)=\Tr'(A)+\Tr'(iB)=\Tr'(A)+i\Tr'(B)
=\TrM(A)+i\TrM(B)=\TrM(A+iB)=\TrM(C)$
or $\Tr'(C)=\TrM(C)$.

\chapterhead{III}{The exact form of $\TrM(A)$ (for $\alpha\geq1$)}%
  {Chapter III. The exact form of Tr M(A) (for α ≥ 1)}

\SectionStart{3.1}{3.1. The spectral form ∫₀¹ φ(λ) dE(λ) with D M(E(λ)) = λ}
3.1. We again assume $M$ in case $\mathrm{II}_1$. Some lemmas about families of
projections and spectral forms in $M$ are needed. Of these the two last ones
have some interest of their own.

\begin{statement}{lemma:3.1.1}{Lemma 3.1.1}
For any two projections $E$, $F\in M$ with $E\leq F$, and any
$\alpha\geq\DM(E)$, $\leq\DM(F)$, a projection $G\in M$ with
$E\leq G\leq F$, $\DM(G)=\alpha$ exists.
\end{statement}

$F-E$ is a projection, and
$0\leq\alpha-\DM(E)\leq\DM(F)-\DM(E)=\DM(F-E)\leq1$.
Thus a projection $G'\in M$ with $\DM(G')=\alpha-\DM(E)$ exists (because $M$ is
in case $\mathrm{II}_1$). So $\DM(G')\leq\DM(F-E)$ and therefore a projection
$G''\in M$ with $G''\leq F-E$, $\DM(G'')=\DM(G')=\alpha-\DM(E)$ exists.
$G''$ is thus orthogonal to $E$, and therefore $G''+E$ is a projection and
$\DM(G''+E)=\DM(G'')+\DM(E)=\alpha$. Besides, $G''+E\geq E$ and
$G''+E\leq(F-E)+E=F$. So $G=G''+E$ has the desired properties.

\begin{statement}{lemma:3.1.2}{Lemma 3.1.2}
For any two projections $E$, $F\in M$ with $E\leq F$ a family of projections
$G(\alpha)\in M$ defined for all $\alpha$ with $\DM(E)\leq\alpha\leq\DM(F)$ exists,
which possesses the following properties:
\begin{enumerate}[label=(\roman*)]
 \item $G(\alpha)=E$ or $F$ if $\alpha=\DM(E)$ or $\DM(F)$ respectively.
 \item $\alpha\leq\beta$ implies $G(\alpha)\leq G(\beta)$.
 \item $\DM(G(\alpha))=\alpha$.
\end{enumerate}
\end{statement}

Choose a sequence $\rho_1,\rho_2,\cdots$ which lies and is dense in the interval
$\DM(E)\leq\alpha\leq\DM(F)$, with $\rho_1=\DM(E)$, $\rho_2=\DM(F)$.

We will now define a sequence of projections $G(\rho_1),G(\rho_2),\cdots$, $\in M$, so
that (i)--(iii) hold for $\alpha=\rho_1,\rho_2,\cdots$. Put first $G(\rho_1)=E$,
$G(\rho_2)=F$; then (i)--(iii) hold for $\alpha=\rho_1$, $\rho_2$. Assume now that for a
$j=3,4,\cdots$ the $G(\rho_1),\cdots,G(\rho_{j-1})$ are already defined, so that
% Editorial note: the source uses the singular verb ``holds'' here; retained verbatim.
(i)--(iii) holds for $\alpha=\rho_1,\cdots,\rho_{j-1}$, we will now define
$G(\rho_j)$ without violating (i)--(iii).

Consider the $\rho_i$, $i=1,\cdots,j-1$, with $\rho_i\leq\rho_j$ (they exist:
$\rho_1=\DM(E)\leq\rho_j$); let $\rho_{i'}$ be the greatest one. Consider the $\rho_i$,
$i=1,\cdots,j-1$, with $\rho_i\geq\rho_j$ (they exist:
$\rho_2=\DM(F)\geq\rho_j$); let $\rho_{i''}$ be the smallest one. Thus
$\rho_{i'}\leq\rho_{i''}$, and so $G(\rho_{i'})\leq G(\rho_{i''})$. Besides
$\DM(G(\rho_{i'}))=\rho_{i'}\leq\rho_j\leq\rho_{i''}=\DM(G(\rho_{i''}))$.
So \StatementRef{lemma:3.1.1}{Lemma 3.1.1} can be applied to
$E=G(\rho_{i'})$, $F=G(\rho_{i''})$, $\alpha=\rho_j$, and we define $G(\rho_j)=G$.

Thus $G(\rho_{i'})\leq G(\rho_j)\leq G(\rho_{i''})$, therefore
$\rho_i\leq\rho_j$, $i=1,\cdots,j-1$, implies
$\rho_i\leq\rho_{i'}$, $G(\rho_i)\leq G(\rho_{i'})\leq G(\rho_j)$ and
$\rho_i\geq\rho_j$, $i=1,\cdots,j-1$, implies
$\rho_i\geq\rho_{i''}$, $G(\rho_i)\geq G(\rho_{i''})\geq G(\rho_j)$.
Besides $\DM(G(\rho_j))=\rho_j$. So (i)--(iii) hold for
$\alpha=\rho_1,\cdots,\rho_{j-1},\rho_j$, too.

Thus all $G(\rho_1),G(\rho_2),\cdots$ are defined. Owing to (ii),
$\lim_{\rho_i\to\alpha,\,\rho_i\geq\alpha}G(\rho_i)$ exists for all $\alpha$ with
$\DM(E)=\rho_1\leq\alpha\leq\rho_2=\DM(F)$. If $\alpha$ is equal to a $\rho_j$, then this
limit is meant to denote the value at $\rho_j$. So we can extend the definition of
$G(\alpha)$ to all above $\alpha$ by defining
\[
 G(\alpha)=\lim_{\rho_i\to\alpha,\,\rho_i\geq\alpha}G(\rho_i).
\]

Now (i) remains true, and (ii), (iii) extend by continuity to all above $\alpha$.

\begin{statement}{lemma:3.1.3}{Lemma 3.1.3}
Let $A\in M$ be a definite operator with bound $\leq1$. Then there
exists a monotonous non-decreasing and right semi-continuous function $\phi(\alpha)$ for
$0\leq\alpha\leq1$ with $0\leq\phi(\alpha)\leq1$ and a resolution of identity $E(\alpha)\in M$
(cf. \ROcite{16}, p.~92, or \RO{}, Definition 15.1.1) with the following properties:
\begin{enumerate}[label=(\roman*)]
 \item $E(0)=0$, $E(1)=1$,
 \item $\DM(E(\alpha))=\alpha$,
 \item $(Af,g)=\int_0^1\phi(\alpha)\,d(E(\alpha)f,g)$ for any $f$, $g$, or symbolically,
 \[
   A=\int_0^1\phi(\alpha)\,dE(\alpha).
 \]
\end{enumerate}
\end{statement}

Let $F(\lambda)$ be the resolution of identity corresponding to $A$. As $A$ is definite,
so $F(0)=0$; as the bound of $A$ is $\leq1$, so $F(1)=1$.

Consider now the two following functions:
\[
 \psi(\lambda)=\DM(F(\lambda)),
 \qquad
 \phi(\alpha)=\operatorname*{l.u.b.}_{\substack{0\leq\lambda\leq1\\ \psi(\lambda)\leq\alpha}}\lambda.
\]
$\psi(\lambda)$ is monotonous non-decreasing and right semi-continuous in
$0\leq\lambda\leq1$, because $F(\lambda)$ is a resolution of identity, and besides
$\psi(0)=0$, $\psi(1)=1$. Thus $\phi(\alpha)$ is monotonous non-decreasing and right
semi-continuous in $0\leq\alpha\leq1$, and its values are all in $0\leq\lambda\leq1$.
One verifies these facts, as well as those which follow, easily; besides, their
discussion may be found in \ROcite{19} (numbers in parentheses refer to the bibliography
% Editorial note: the source's singular verb ``corresponds'' with a compound subject is retained verbatim.
in \RO{}, pp.~125--126) on p.~193. (Our $\psi(\lambda)$, $\phi(\alpha)$ corresponds to the
$\phi(a)$, $\psi(b)$ there, thus they are both ``$m$-functions,'' and $\phi(\alpha)$ is the
``measure function'' of $\psi(\lambda)$.) The relation between $\phi(\alpha)$ and
$\psi(\lambda)$ is symmetric (cf. loc. cit., $\psi(\lambda)$ is the ``measure function''
of $\phi(\alpha)$, we must define for the empty set in $0\leq\lambda\leq1$ the l.u.b.~0);
that is,
\[
 \psi(\lambda)=\operatorname*{l.u.b.}_{0\leq\alpha\leq1}(\phi(\alpha)\leq\lambda).
\]
Finally (cf. loc. cit.)
\EquationTarget{star}
\begin{equation*}
 \psi(\phi(\psi(\lambda)))=\psi(\lambda). \tag{$*$}
\end{equation*}
That is, $\DM(F(\phi(\psi(\lambda))))=\DM(F(\lambda))$. But
$F(\phi(\psi(\lambda)))\geq F(\lambda)$, and therefore this implies
\EquationTarget{starstar}
\begin{equation*}
 F(\phi(\psi(\lambda)))=F(\lambda). \tag{$*\,*$}
\end{equation*}

If $\phi(\alpha)<1$, then choose $\mu$ with $\phi(\alpha)<\mu\leq1$, and let
$\mu\to\phi(\alpha)$. Under these conditions the definition of $\phi(\alpha)$
implies $\DM(F(\mu))>\alpha$, so $\DM(F(\phi(\alpha)))\geq\alpha$. For $\phi(\alpha)=1$
this holds, too: $\DM(F(\phi(\alpha)))=\DM(F(1))=\DM(1)=1\geq\alpha$. On the other hand,
if $\phi(\alpha)>0$, then a sequence of $\mu$ with $0\leq\mu\leq\phi(\alpha)$,
$\mu\to\phi(\alpha)$ and $\DM(F(\mu))\leq\alpha$ exists, so
$\DM(F(\phi(\alpha)-0))\leq\alpha$.
% Editorial note: source prints phi(0)=0 here (rather than phi(alpha)=0); retained verbatim.
This holds for $\phi(0)=0$, too, if we identify
$F(0-0)$ with $F(0)=0$: $\DM(F(0-0))=\DM(0)=0\leq\alpha$. So we have
\EquationTarget{hash}
\begin{equation*}
 \DM(F(\phi(\alpha)-0))\leq\alpha\leq\DM(F(\phi(\alpha))). \tag{$\#$}
\end{equation*}



Form now for every $0\leq\lambda\leq1$, $E=F(\lambda-0)$,
$F=F(\lambda)$ and apply \StatementRef{lemma:3.1.2}{Lemma~3.1.2}.
A family of projections $G(\alpha)\in M$,
$D_M(F(\lambda-0))\leq\alpha\leq D_M(F(\lambda))$, results. Choose for
every $0\leq\lambda\leq1$ such a family, and hold it
fixed: $G(\lambda,\alpha)$. (If $F(\lambda-0)=F(\lambda)$, then necessarily
$\alpha=D_M(F(\lambda))$, and $G(\lambda,\alpha)=F(\lambda)$. So only the
$\lambda$ with $F(\lambda-0)\ne F(\lambda)$, the point-proper values of $A$,
are important. Their set is finite or enumerably infinite.)

Now define for every $0\leq\alpha\leq1$
\[
 \EquationTarget{3.1.section}(\S)\qquad
 E(\alpha)=G(\phi(\alpha),\alpha)
\]
(this expression has a meaning owing to \EquationRef{hash}{(\#)}).

We have now by definition $D_M(E(\alpha))=\alpha$, so condition (ii) is
satisfied. For $\alpha=0,1$, this gives $E(0)=0$, $E(1)=1$, so condition (i)
is satisfied, too. If $\alpha\leq\beta$, then
$\phi(\alpha)\leq\phi(\beta)$. The relation
$\phi(\alpha)=\phi(\beta)$ gives
\[
 E(\alpha)=G(\phi(\alpha),\alpha)
 =G(\phi(\beta),\alpha)
 \leq G(\phi(\beta),\beta)=E(\beta),
\]
while $\phi(\alpha)<\phi(\beta)$ gives
\[
 E(\alpha)=G(\phi(\alpha),\alpha)
 \leq F(\phi(\alpha))
 \leq F(\phi(\beta)-0)
 \leq G(\phi(\beta),\beta)=E(\beta).
\]
So we see that
\[
 \EquationTarget{3.1.bullet}(\mathord{\cdot})\qquad
 \alpha\leq\beta\quad\hbox{implies}\quad E(\alpha)\leq E(\beta).
\]
Furthermore $\beta>\alpha$ gives
\[
 D_M(E(\beta)-E(\alpha))
 =D_M(E(\beta))-D_M(E(\alpha))=\beta-\alpha,
\]
so
$\lim_{\substack{\beta\to\alpha\\\beta>\alpha}}
D_M(E(\beta)-E(\alpha))=0$. But
\[
 \lim_{\substack{\beta\to\alpha\\\beta>\alpha}}
 D_M(E(\beta)-E(\alpha))
 =D_M\left(\lim_{\substack{\beta\to\alpha\\\beta>\alpha}}E(\beta)
 -E(\alpha)\right)
 =D_M(E(\alpha+0)-E(\alpha)),
\]
so $D_M(E(\alpha+0)-E(\alpha))=0$, $E(\alpha+0)=E(\alpha)$, or
\[
 \EquationTarget{3.1.doublebullet}(\mathord{\cdot}\mathord{\cdot})\qquad
 \lim_{\substack{\beta\to\alpha\\\beta>\alpha}}E(\beta)=E(\alpha).
\]

(i), \EquationRef{3.1.bullet}{$(\mathord{\cdot})$},
\EquationRef{3.1.doublebullet}{$(\mathord{\cdot}\mathord{\cdot})$} state
that $E(\alpha)$ is a resolution of unity.

It remains for us to prove (iii).

Consider the expression
$\int_0^1\phi(\alpha)\,d(\lVert E(\alpha)f\rVert^2)$.
$\phi(\alpha)$ and $\lVert E(\alpha)f\rVert^2$ are both monotonous
non-decreasing functions of $\alpha$; if $\alpha$ increases from 0 to 1,
then these functions increase from $\phi(0)$ to $\phi(1)=1$ and from 0 to
$\lVert f\rVert^2$ respectively. Thus our integral is, by partial
integration, equal to
$\lVert f\rVert^2-\int_0^1\lVert E(\alpha)f\rVert^2\,d\phi(\alpha)$.
We may now introduce a new variable of integration: $\lambda=\phi(\alpha)$.
Then we have (cf. loc. cit., p.~198)
$\int_0^1\lVert E(\alpha)f\rVert^2\,d\phi(\alpha)
=\int_0^1\lVert E(\psi(\lambda))f\rVert^2\,d\lambda$.
But by \EquationRef{3.1.section}{(\S)},
$E(\psi(\lambda))=G(\phi(\psi(\lambda)),\psi(\lambda))$, and by
\EquationRef{star}{(*)},
$\psi(\phi(\psi(\lambda)))=D_M(F(\phi(\psi(\lambda))))$. So the definition
of $G(\mu,\alpha)$ gives
$G(\phi(\psi(\lambda)),\psi(\lambda))
=F(\phi(\psi(\lambda)))$.
By \EquationRef{starstar}{(**)} this is $F(\lambda)$, so we have
$E(\psi(\lambda))=F(\lambda)$. Thus our original expression is equal to
$\lVert f\rVert^2-\int_0^1\lVert F(\lambda)f\rVert^2\,d\lambda$,
and this becomes by another partial integration
$\int_0^1\lambda\,d\lVert F(\lambda)f\rVert^2$, that is, $(Af,f)$.
Thus we have proved:
\[
 (Af,f)=\int_0^1\phi(\alpha)\,d\lVert E(\alpha)f\rVert^2.
\]

Replacing herein $f$ by $f\pm g/2$, and subtracting, gives the real part of
(iii). Now replacing $f,g$ by $if,g$ gives the imaginary part of (iii).
Thus the proof is completed.

Finally we need the following lemma.

\begin{statement}{lemma:3.1.4}{Lemma 3.1.4}
Let $E(\alpha)\in M$ be a resolution of unity in $a\leq\alpha\leq b$,
$E(a)=0$, $E(b)=1$. Then a bounded Hermitian operator $B\in M$ with
$(Bf,g)=\int_a^b\lambda\,d(E(\lambda)f,g)$,
symbolically $B=\int_a^b\lambda\,dE(\lambda)$, exists. For every bounded
Baire function $\psi(\alpha)$ an operator $C=\psi(B)\in M$ with
$(Cf,g)=\int_a^b\psi(\lambda)\,d(E(\lambda)f,g)$,
symbolically $C=\int_a^b\psi(\lambda)\,dE(\lambda)$ exists. If
$0\leq\psi(\lambda)\leq1$, then $C$ is Hermitian, definite, and of bound 1.
For all $\psi(\lambda)$
\[
 \TrM(C)=\int_a^b\psi(\lambda)\,dD_M(E(\lambda)).
\]
\end{statement}

The existence of the bounded $B$ is a well known fact about spectral forms,
as all $E(\lambda)\in M$ so $B\in M$ (see \ROcite{18}, p.~389,
Theorem~1). $C=\psi(B)$ exists and is bounded (\ROcite{19}, p.~205), and
$B\in M$ implies $C\in M$ (\ROcite{19}, p.~213, Theorem~6).
$0\leq\psi(\alpha)\leq1$ implies the Hermitian, definite character of $C$
and $1-C$ (\ROcite{19}, p.~205) ($C$ Hermitian: by Property b; $C$ and
$1-C$ definite: by Property e), with $F(x)=\psi(x)$ or $1-\psi(x)$ and
$G(x)=H(x)=(F(x))^{1/2}$. Thus the bound of $C$ is $\leq1$.
(Alternatively \ROcite{16}, p.~113, Theorem~4*, could be used.)

It remains to prove the final formula for the trace. If it holds for
$\psi_1(\lambda),\psi_2(\lambda),\cdots$, and the $\psi_n(\lambda)$,
$n=1,2,\cdots$ are uniformly bounded and everywhere convergent in
$a\leq\lambda\leq b$, then it holds for
$\psi(\lambda)=\lim_{n\to\infty}\psi_n(\lambda)$ too:
$\psi(B)=\mathop{\mathrm{strong}\lim}_{n\to\infty}\psi_n(B)$,
the $\psi_n(B)$ being uniformly bounded, by \ROcite{19}, p.~205,
Property h), and $\TrM$ is continuous for strong (even for
weak) convergence by \StatementRef{property:III}{Property III},
\EquationRef{property:III:starstar}{$(*\,*)$}. Thus our relation holds for all
bounded Baire functions $\psi(\lambda)$ if it holds for all continuous
functions $\psi(\lambda)$. And it holds for these, if it holds for all
intervalwise constant functions $\psi(\lambda)$. But these are linear
% Editorial note: source has no punctuation between d and ``=0 otherwise''; retained verbatim.
aggregates of functions $\psi_{c,d}(\lambda)=1$ for $c<\lambda\leq d$
$=0$ otherwise, where $a\leq c\leq d\leq b$. So it suffices to consider
these. Now clearly
\begin{align*}
 \psi_{c,d}(B)&=E(d)-E(c),\\
 \TrM(\psi_{c,d}(B))
 &=D_M(E(d)-E(c))=D_M(E(d))-D_M(E(c)),
\end{align*}
and
\[
 \int_a^b\psi_{c,d}(\alpha)\,dD_M(E(\alpha))
 =D_M(E(d))-D_M(E(c)),
\]
completing the proof.

\SectionStart{3.2}{3.2. The maximum problems (A) and (B)}3.2. We return now to the normalization of
\SectionRef{1.1}, and assume that an $M$ in case $\mathrm{II}_1$
is given with $\alpha\geq1$.

Consider a $g$ and $K$ which are related as in
\StatementRef{theorem:I}{Theorem~I}. Our objective is to prove
\StatementRef{theorem:II}{Theorem~II} below, but we begin by considering
these two maximum problems.

\indent\ignorespaces
(A) For a given $\lambda$ with $0\leq\lambda\leq1$ consider all projections
$F\in M$ with $D_M(F)=\lambda$, and the corresponding values of $(Fg,g)$.
Prove that these $(Fg,g)$ possess a maximum $m_a$ which they assume for a
certain $F=F_0$ from the above class.

\indent\ignorespaces
(B) For a given $\lambda$ with $0\leq\lambda\leq1$ consider all definite
operators $B\in M$ of bound $\leq1$ with $\TrM(B)=\lambda$,
and the corresponding values of $(Bg,g)$. Prove that these $(Bg,g)$ possess
a maximum $m_b$ which they assume for a certain $B=B_0$ from the above
class.

\begin{statement}{lemma:3.2.1}{Lemma 3.2.1}
Problem (B) possesses a solution $B=B_0$, the maximum $m_b$ fulfills
$\lambda K^{-1}\leq m_b\leq\lambda K$.
\end{statement}

For the $B$ of the class described in Problem (B) we have by
\StatementRef{theorem:I}{Theorem~I},
$\lambda K^{-1}\leq(Bg,g)\leq\lambda K$.
So these $(Bg,g)$ possess a l.u.b. $m_b'$ and
$\lambda K^{-1}\leq m_b'\leq\lambda K$.

Select now a sequence $B_1,B_2,\cdots$ from this class, so that
$\lim_{n\to\infty}(B_ng,g)=m_b'$. As the $B_1,B_2,\cdots$ are uniformly
bounded, there exists a subsequence $B_{n_1},B_{n_2},\cdots$
($n_1<n_2<\cdots$) such that
$B_0=\mathop{\mathrm{weak}\lim}_{i\to\infty}B_{n_i}$
exists. As all $B_{n_i}$ are in $M$, definite, and of bound $\leq1$, the
same is true for $B_0$. As all $\TrM(B_{n_i})=\lambda$, so
\StatementRef{property:II}{Property II} or
\StatementRef{property:III}{III},
\EquationRef{property:III:starstar}{$(*\,*)$}, gives $\TrM(B_0)=\lambda$.
Finally
$(B_0g,g)=\lim_{i\to\infty}(B_{n_i}g,g)
=\lim_{n\to\infty}(B_ng,g)=m_b'$.

Thus $B_0$ belongs to our class, and $(B_0g,g)=m_b'$. Therefore the l.u.b.
$m_b'$ is a maximum: $m_b'=m_b$, and
$\lambda K^{-1}\leq m_b\leq\lambda K$. So the proof is completed.

\begin{statement}{lemma:3.2.2}{Lemma 3.2.2}
The $B_0$ of \StatementRef{lemma:3.2.1}{Lemma~3.2.1} can be chosen as a
projection.
\end{statement}

Consider the $B_0$ of \StatementRef{lemma:3.2.1}{Lemma~3.2.1}. By
\StatementRef{lemma:3.1.3}{Lemmas~3.1.3} and
\StatementRef{lemma:3.1.4}{3.1.4} we have $B_0=\phi(B)$, where
$B=\int_0^1\alpha\,dE(\alpha)$ (symbolically),
$E(\alpha)$ being a resolution of unity with $E(\alpha)\in M$,
$D_M(E(\alpha))=\alpha$, and $\phi(\alpha)$ a monotonous non-decreasing and
right semi-continuous function.

Consider any Baire function $\psi(\alpha)$, $0\leq\alpha\leq1$, with
$0\leq\psi(\alpha)\leq1$. Form
$\psi(B)=\int_0^1\psi(\alpha)\,dE(\alpha)$ (symbolically),
using \StatementRef{lemma:3.1.4}{Lemma~3.1.4}. Then $\psi(B)$ is in $M$,
definite, and of bound $\leq1$, and
\[
 \TrM(\psi(B))
 =\int_0^1\psi(\alpha)\,dD_M(E(\alpha))
 =\int_0^1\psi(\alpha)\,d\alpha.
\]

Thus $\psi(B)$ belongs to the class of Problem (B) if and only if
$\int_0^1\psi(\alpha)\,d\alpha=\lambda$. Besides
$(\psi(B)g,g)=\int_0^1\psi(\alpha)\,d(\lVert E(\alpha)g\rVert^2)$.

Apply this to $\psi(\alpha)=\phi(\alpha)$,
$\psi(B)=\phi(B)=B_0$. We get
$\int_0^1\phi(\alpha)\,d\alpha=\lambda$,
$\int_0^1\phi(\alpha)\,d(\lVert E(\alpha)g\rVert^2)=m_b$.
And the maximum property of $B_0$ gives the following. For every Baire
function $\psi(\alpha)$, $0\leq\alpha\leq1$, with
$0\leq\psi(\alpha)\leq1$ the equation
$\int_0^1\psi(\alpha)\,d\alpha=\lambda$ implies
$\int_0^1\psi(\alpha)\,d\lVert E(\alpha)g\rVert^2\leq m_b$.

If we had $\phi(1-\lambda)=0$, then the monotony of $\phi(\alpha)$ would
give $\phi(\alpha)=0$ for $0\leq\alpha\leq1-\lambda$, and the right
semi-continuity $\phi(\alpha)\leq\tfrac12$ for
$1-\lambda<\alpha\leq1-\lambda+\epsilon'$ for a suitable $\epsilon'>0$.
Besides $\phi(\alpha)$ is always $\leq1$. Thus
\[
 \int_0^1\phi(\alpha)\,d\alpha
 \leq0(1-\lambda)+\tfrac12\epsilon'+(\lambda-\epsilon')1
 =\lambda-\tfrac12\epsilon'<\lambda,
\]
contradicting $\int_0^1\phi(\alpha)\,d\alpha=\lambda$. Therefore
$\phi(1-\lambda)>0$.

If we had $\phi(1-\lambda-0)=1$, then the monotony of $\phi(\alpha)$ would
give $\phi(\alpha)=1$ for $1-\lambda\leq\alpha\leq1$ and the definition of
$\phi(1-\lambda-0)$, $\phi(\alpha)\geq\tfrac12$ for
$1-\lambda-\epsilon'\leq\alpha<1-\lambda$, for a suitable $\epsilon'>0$.
Besides $\phi(\alpha)$ is always $\geq0$. Thus
\[
 \int_0^1\phi(\alpha)\,d\alpha
 \geq(1-\lambda-\epsilon')\cdot0+\frac{\epsilon'}2+\lambda\cdot1
 =\lambda+\frac{\epsilon'}2>\lambda,
\]
contradicting $\int_0^1\phi(\alpha)\,d\alpha=\lambda$. Therefore
$\phi(1-\lambda-0)<1$.

So we can choose a $\delta$, $0<\delta<1$, with
$\phi(1-\lambda)>\delta$, $\phi(1-\lambda-0)<1-\delta$.
Thus $1-\lambda\leq\alpha\leq1$ implies
$(\phi(\alpha)-\delta)/(1-\delta)
\geq(\phi(1-\lambda)-\delta)/(1-\delta)\geq0$ and
$\leq(1-\delta)/(1-\delta)=1$,
% Editorial note: the source uses \leq at \alpha=1-\lambda in both adjoining
% cases (and again in the definition of \phi_1), so the cases overlap;
% retained verbatim.
and $0\leq\alpha\leq1-\lambda$ implies
$\phi(\alpha)/(1-\delta)\geq0$ and $\leq(1-\delta)/(1-\delta)=1$.
Now put $\phi_1(\alpha)=1$ for $1-\lambda\leq\alpha\leq1$ and
$\phi_1(\alpha)=0$ for $0\leq\alpha\leq1-\lambda$. Then we have
$0\leq\phi_1(\alpha)\leq1$ for all $0\leq\alpha\leq1$; and also, for
$\phi_2(\alpha)=(\phi(\alpha)-\delta\phi_1(\alpha))/(1-\delta)$,
$0\leq\phi_2(\alpha)\leq1$ for all $0\leq\alpha\leq1$ (we proved this
above separately for $1-\lambda\leq\alpha\leq1$ and for
$0\leq\alpha\leq1-\lambda$). Besides
\[
 \EquationTarget{3.2.star}(*)\qquad
 \phi(\alpha)=\delta\phi_1(\alpha)+(1-\delta)\phi_2(\alpha).
\]

Now clearly $\int_0^1\phi_1(\alpha)\,d\alpha=\lambda$. This and
$\int_0^1\phi(\alpha)\,d\alpha=\lambda$ and
\EquationRef{3.2.star}{(*)} give
$\int_0^1\phi_2(\alpha)\,d\alpha=\lambda$. Therefore
$\int_0^1\phi_1(\alpha)\,d\lVert E(\alpha)g\rVert^2\leq m_b$,
$\int_0^1\phi_2(\alpha)\,d\lVert E(\alpha)g\rVert^2\leq m_b$.
But \EquationRef{3.2.star}{(*)} gives
\begin{align*}
 &\delta\int_0^1\phi_1(\alpha)\,d\lVert E(\alpha)g\rVert^2
 +(1-\delta)\int_0^1\phi_2(\alpha)\,d\lVert E(\alpha)g\rVert^2\\
 &\qquad=
 \int_0^1\phi(\alpha)\,d\lVert E(\alpha)g\rVert^2=m_b.
\end{align*}
Therefore necessarily
$\int_0^1\phi_1(\alpha)\,d\lVert E(\alpha)g\rVert^2=m_b$.
In other words, for $B_1=\phi_1(B)$ which belongs to the class of Problem
(B), we have $(B_1g,g)=m_b$.

% Editorial note: the source prints B_1=\phi_1(B)=1-E(\lambda), although the
% preceding cutoff defining \phi_1 is 1-\lambda; retained verbatim.
Thus $B_1=\phi_1(B)$ is a solution of Problem (B), too. But one verifies
easily that $B_1=\phi_1(B)=1-E(\lambda)$, and so $B_1$ is a projection.

Therefore replacement of $B_0$ by $B_1$ completes the proof.

\begin{statement}{lemma:3.2.3}{Lemma 3.2.3}
Problems (A) and (B) possess a common solution $E_0=B_0$, and for the
maxima we have $m_a=m_b$.
\end{statement}

The class of the $E$ in Problem (A) is obviously a subclass of the $B$ in
Problem (B): It consists of all projections of the latter class. (This is so,
because for projections $E$, $D_M(E)=\TrM(E)$
(cf. \RO{}, Lemma~15.3.1).) Now
\StatementRef{lemma:3.2.2}{Lemma~3.2.2} states, that a solution of Problem
(B) exists, which is a projection and thus belongs to the class of Problem
(A). Therefore it is a solution of Problem (A) too, and $m_a=m_b$.

Problems (A) and (B) can be modified, when a projection
$G=P_{\mathfrak N}\in M$ is given, in the following sense: Replace
$\mathfrak H,M$ by $\mathfrak N,M_{(\mathfrak N)}$
(cf. \RO{}, \S11.3). Then $D_M(E)$ must be replaced by
$D_M(E)/D_M(G)$, in order to conserve the standard normalization and so
the normalization of \SectionRef{1.1} requires replacing of
% Editorial note: source reads the slightly awkward ``This alpha''; retained verbatim.
$D_{M'}(E')$ by $D_{M'}(E')/D_M(G)$. This $\alpha$ is replaced by
$\alpha/D_M(G)$ and so $\alpha\geq1$ remains true.
\StatementRef{lemma:3.2.3}{Lemma~3.2.3} is modified, in so far as
$\lambda$ is replaced by $\lambda/D_M(G)$ and $M$ in Problems (A) and (B)
is to be replaced by $M_{(\mathfrak N)}$. This means that projections
$F\in M$ must be $\leq G$, and operators $B\in M$ must fulfill
$BG=GB=B$.

Combining this and the corresponding changes in
\StatementRef{lemma:3.2.1}{Lemma~3.2.1} (observe that
$(B_0g,g)=(E_0g,g)=\lVert E_0g\rVert^2$), we have

\begin{statement}{corollary:3.2.3}{Corollary}
Let a projection $G\in M$ be given. Replace Problems (A) and (B) by
Problems $(\mathrm A_G)$ and $(\mathrm B_G)$, which arise by imposing these
further restrictions: In $(\mathrm A_G)$ the projections $F\in M$ are
$F\leq G$. In $(\mathrm B_G)$ the operators $B\in M$ fulfill $BG=GB=B$.
Assume $0\leq\lambda\leq D_M(G)$. Then Problems $(\mathrm A_G)$ and
$(\mathrm B_G)$ possess a common solution $E_0=B_0$, and we have for the
maxima
$K^{-1}\lVert E_0g\rVert^2\leq m_a=m_b\leq K\lVert E_0g\rVert^2$.
\end{statement}

\SectionStart{3.3}{3.3. Reducibility of a g ∈ H by an E ∈ M}3.3. We hold $g$ fixed, and make the following

\begin{statement}{definition:3.3.1}{Definition 3.3.1}
Let $E,G$ be two projections in $M$, $E\leq G$. We say that $E$ reduces
$g$ with respect to $G$ if for every $A\in M$ with $AG=GA=A$
\[
 (AEg,g)=(EAg,g).
\]

If $G$ may be taken $=1$, we say that $E$ reduces $g$.
\end{statement}

\begin{statement}{lemma:3.3.1}{Lemma 3.3.1}
Let $G$ be a projection in $M$, $E_0$ a solution of Problem
$(\mathrm A_G)$ (cf. the \StatementRef{corollary:3.2.3}{corollary} to
\StatementRef{lemma:3.2.3}{Lemma~3.2.3}). Then $E_0$ reduces $g$ with
respect to $G$.
\end{statement}

If $U$ is unitary, in $M$, and commutes with $G$, then $U^{-1}E_0U$ is a
projection in $M$, it is $\leq U^{-1}GU=G$, and
$D_M(U^{-1}E_0U)=D_M(E_0)=\lambda$. So
$(U^{-1}E_0Ug,g)\leq(E_0g,g)$.
Now
$(U^{-1}E_0Ug,g)=(U^*E_0Ug,g)=(E_0Ug,Ug)$.
So $(E_0g,g)-(E_0Ug,Ug)\geq0$.

% Editorial note: source reads \epsilon\gtrless0 and
% \psi_\epsilon(\alpha)=-\alpha^2/(1-i\epsilon\alpha); retained verbatim.
Let $A$ be Hermitian, in $M$, and commute with $G$. Put for some
$\epsilon\gtrless0$, $\phi_\epsilon(\alpha)=(1+i\epsilon\alpha)/(1-i\epsilon\alpha)$.
As $|\phi_\epsilon(\alpha)|=1$, so $U_\epsilon=\phi_\epsilon(A)$ is
unitary, and it is in $M$ and commutes with $G$ along with $A$. Further
$\phi_\epsilon(\alpha)=1+2i\epsilon\alpha+\epsilon^2\psi_\epsilon(\alpha)$,
where $\psi_\epsilon(\alpha)=-\alpha^2/(1-i\epsilon\alpha)$.
So $|\psi_\epsilon(\alpha)|\leq D^2$ if $|\alpha|\leq D$, and therefore
$\opnorm{\psi_\epsilon(A)}\leq\opnorm{A}^2$
(cf., for instance, \ROcite{16}, p.~113, Theorem~4*). Now
\begin{align*}
 0&\leq(E_0g,g)-(E_0U_\epsilon g,U_\epsilon g)=(E_0g,g)\\
 &\quad-(E_0(1+2i\epsilon A+\epsilon^2\psi_\epsilon(A))g,
 (1+2i\epsilon A+\epsilon^2\psi_\epsilon(A))g)\\
 &=2\epsilon(i(AE_0-E_0A)g,g)+O(\epsilon^2).
\end{align*}
So we have
\[
 2\epsilon i((AE_0-E_0A)g,g)+O(\epsilon^2)\geq0.
\]

As $\epsilon\gtrless0$, this necessitates
$((AE_0-E_0A)g,g)=0$, that is,
\[
 (AE_0g,g)=(E_0Ag,g).
\]
This equation extends from the Hermitian $A\in M$ to all $A\in M$, which
commute with $G$. Put for such an $A$,
$A_1=\tfrac12(A+A^*)$, $A_2=-\tfrac12i(A-A^*)$,
then it holds for $A_1,A_2$ and so for $A=A_1+iA_2$. Thus we have
established that $E_0$ reduces $g$ with respect to $G$.

\begin{statement}{lemma:3.3.2}{Lemma 3.3.2}
If $E$ reduces $g$ relatively to $G$ and $G$ reduces $g$, then $E$ reduces
$g$.
\end{statement}

We must show that for every $A\in M$, $(AEg,g)=(EAg,g)$. Now
\begin{align*}
 (EAg,g)&=(GEAg,g)=(G(EA)g,g)=((EA)Gg,g)=(EGAGg,g)\\
 &=(GAGEg,g)=(GAEg,g)=((AE)Gg,g)=(AEg,g).
\end{align*}

\begin{statement}{lemma:3.3.3}{Lemma 3.3.3}
If $\{E_i\}$ is a sequence of projections $E_i\in M$ each of which reduce
$g$ and $E_i\cdot E_j=0$ if $i\ne j$, then $\sum_{i=1}^\infty E_i$ reduces $g$.
If $E_1$ and $E_2$ reduce $g$ and $E_1\geq E_2$ then $E_1-E_2$ reduces
$g$.
\end{statement}

This is immediate for $A\in M$ implies $A$ is bounded.

\begin{statement}{lemma:3.3.4}{Lemma 3.3.4}
Let $\{E_i\}$ be a sequence with the same properties as in
\StatementRef{lemma:3.3.3}{Lemma~3.3.3}. Let
$E=\sum_{i=1}^\infty E_i$. Then for $A\in M$,
\[
 (EAg,g)=(AEg,g)=(EAEg,g)=\sum_{i=1}^\infty(E_iAE_ig,g).
\]
Also if $E$ reduces $g$, $EF=0$, $F\in M$, then
$(EAFg,g)=(FAEg,g)=0$.
\end{statement}

We have
$(EAg,g)=(E(EA)g,g)=(EAEg,g)=(E(AE)g,g)=((AE)Eg,g)=(AEg,g)$.
Hence we have
$(EAg,g)=\left(\left(\sum_{i=1}^\infty E_i\right)Ag,g\right)
=\sum_{i=1}^\infty(E_iAg,g)=\sum_{i=1}^\infty(E_iAE_ig,g)$.
Also
$(EAFg,g)=(E(AF)g,g)=((AF)Eg,g)=(A(FE)g,g)=0
=((EF)Ag,g)=(E(FA)g,g)=((FA)Eg,g)=(FAEg,g)$.

\begin{statement}{lemma:3.3.5}{Lemma 3.3.5}
Let $\{E_i\}$ be a sequence of mutually orthogonal projections. If
$E=\sum_{i=1}^\infty E_i$ and $A$ and $E_i\in M$,
\[
 \TrM(EAE)=\sum_{i=1}^\infty\TrM(E_iAE_i).
\]
\end{statement}

By \StatementRef{property:III}{Property III}, (vi),
\(\TrM(EAE)=\TrM(AE\cdot E)=\TrM(AE)\).
Similarly $\TrM(E_iAE_i)=\TrM(AE_i)$.
By \StatementRef{property:III}{Property III},
\EquationRef{property:III:starstar}{$(*\,*)$}, and (iii)
\(\TrM(AE)=\TrM\left(\sum_{i=1}^\infty AE_i\right)
\allowbreak=\sum_{i=1}^\infty\TrM(AE_i)\).
Thus
$\TrM(EAE)=\sum_{i=1}^\infty\TrM(E_iAE_i)$.
Note that we have also demonstrated the convergence of
$\sum_{i=1}^\infty\TrM(E_iAE_i)$.

\begin{statement}{lemma:3.3.6}{Lemma 3.3.6}
Let $\{E_i\}$ be a sequence of projections each
$E_i\mathrel{\underset p{\geq}}\lambda$
($\mathrel{\underset p{\leq}}\lambda$) for $g$. Furthermore let us suppose
that each $E_i$ reduces $g$ and $E_i\cdot E_j=0$ if $i\ne j$. Then
\(E=\sum_{i=1}^\infty E_i\mathrel{\underset p{\geq}}\lambda
(\mathrel{\underset p{\leq}}\lambda)\).
\end{statement}

Let $A$ be positive definite, in $M$ and such that $EA=AE=A$. Then by
\StatementRef{lemma:3.3.4}{Lemma~3.3.4}
\[
 (Ag,g)=(EAg,g)=\sum_{i=1}^\infty(E_iAE_ig,g).
\]
Now $E_iAE_i$ is positive definite and is unchanged by left- or
right-multiplication with $E_i$. Hence
\StatementRef{lemma:1.4.2}{Lemma~1.4.2} yields
\((E_iAE_ig,g)\geq\lambda\TrM(E_iAE_i)\).
% Editorial note: source cites the apparently nonexistent Lemma 3.2.6;
% the printed citation is retained, while its link points to Lemma 3.3.5,
% whose trace decomposition is used here.
Thus by \StatementRef{lemma:3.3.5}{Lemma~3.2.6}
\[
 (Ag,g)=\sum_{i=1}^\infty(E_iAE_ig,g)
 \geq\lambda\sum_{i=1}^\infty\TrM(E_iAE_i)
 =\lambda\TrM(A).
\]
\StatementRef{lemma:1.4.2}{Lemma~1.4.2} now implies
$E\mathrel{\underset p{\geq}}\lambda$.

\SectionStart{3.4}{3.4. Decomposition of a g from Theorem I}3.4. We still suppose that $g$ is held fixed.

\begin{statement}{lemma:3.4.1}{Lemma 3.4.1}
Let $E\ne0$ be a projection, in $M$, and such that
$a\mathrel{\underset p{\leq}}E\mathrel{\underset p{\leq}}b$. Let a
$\lambda$ with $0<\lambda<D_M(E)$ be given. Then there exists a projection
$E_0\in M$ with the following properties: (i) $E_0\leq E$;
(ii) $D_M(E_0)=\lambda$; (iii) $E_0$ reduces $g$ with respect to $E$; and
(iv) there exists a $\xi_0$ with
\(a\mathrel{\underset p{\leq}}E-E_0
\mathrel{\underset p{\leq}}\xi_0
\mathrel{\underset p{\leq}}E_0
\mathrel{\underset p{\leq}}b\).
\end{statement}

Consider the solution $E_0$ of Problem $(\mathrm A_G)$ with $G=E$, in the
\StatementRef{corollary:3.2.3}{corollary} to
\StatementRef{lemma:3.2.3}{Lemma~3.2.3}. This $E_0$ is a
projection in $M$ and fulfills (i), (ii) by definition, and it fulfills (iii)
by \StatementRef{lemma:3.3.1}{Lemma~3.3.1}. As to (iv), both $E_0$ and
$E-E_0$ are $\mathrel{\underset p{\geq}}a$ and
$\mathrel{\underset p{\leq}}b$ along with $E$, by
\StatementRef{lemma:1.2.3}{Lemma~1.2.3}.

So we must only find a $\xi$ with
$E-E_0\mathrel{\underset p{\leq}}\xi
\mathrel{\underset p{\leq}}E_0$. Let $\Gamma'$ be the set of all $\eta$'s,
for which there exists an $F\leq E_0$ with
$F<\eta$; and let $\Gamma''$ be
the set of all $\eta$'s for which there exists an $F\leq E-E_0$ with
$F>\eta$. $\Gamma'$ and $\Gamma''$ are open
intervals, and clearly every
$\eta>b$ belongs to $\Gamma'$, and every $\eta<a$ belongs to $\Gamma''$.
So there exists either a $\xi_0$ such that every $\eta$ in $\Gamma'$ is
$\geq\xi_0$ and every $\eta$ in $\Gamma''$ is $\leq\xi_0$, or $\Gamma'$
and $\Gamma''$ have a common element $\xi_1$.

If the former holds, then $F\leq E_0$ implies
$F\geq\eta$ for all $\eta<\xi_0$, and so
$F\geq\xi_0$, that is,
$E_0\mathrel{\underset p{\geq}}\xi_0$. Similarly
$E-E_0\mathrel{\underset p{\leq}}\xi_0$. So in this case the rest of (iv)
is proved, too.

If the latter holds, then we have even a $\xi_1-\delta$ in $\Gamma'$ and
$\xi_1+\delta$ in $\Gamma''$ for some suitable $\delta>0$. So an
$F_1\leq E_0$ with $F_1<\xi_1-\delta$ exists, and an
$F_2\leq E-E_0$ with $F_2>\xi_1+\delta$. By
% Editorial note: source cites Lemma 1.2; the argument uses Lemma 1.2.2.
% The printed citation is retained, while its link points to the result used.
\StatementRef{lemma:1.2.2}{Lemma~1.2} there exist two $F_3$ and $F_4$ (both
$\ne0$), with $F_3\leq F_1$, $F_4\leq F_2$ and
\(F_3\mathrel{\underset p{\leq}}\xi_1-\delta\),
\(F_4\mathrel{\underset p{\geq}}\xi_1+\delta\).
We may assume that $D_M(F_3)=D_M(F_4)$, since otherwise we may replace
$F_3,F_4$ by two $F_3'\leq F_3$, $F_4'\leq F_4$ with
\(D_M(F_3')=D_M(F_4')=\min(D_M(F_3),D_M(F_4))\)
(remembering \StatementRef{lemma:1.2.3}{Lemma~1.2.3}).

Now as $F_3\leq E_0$, $F_4\leq E-E_0$, so $E_0-F_3+F_4$ is a projection
$\leq E$, and
\[
 D_M(E_0-F_3+F_4)
 =D_M(E_0)-D_M(F_3)+D_M(F_4)=D_M(E_0)=\lambda,
\]
while
\begin{align*}
 ((E_0-F_3+F_4)g,g)
 &=(E_0g,g)-(F_3g,g)+(F_4g,g)\\
 &\geq m_a-(\xi_1-\delta)D_M(F_3)+(\xi_1+\delta)D_M(F_4)\\
 &\geq m_a-(\xi_1-\delta)D_M(F_3)+(\xi_1+\delta)D_M(F_3)\\
 &=m_a+2\delta D_M(F_3)>m_a,
\end{align*}
contradicting the maximum property of $m_a$. So this case cannot arise.

The proof is therefore completed.

\begin{statement}{lemma:3.4.2}{Lemma 3.4.2}
We can define for all $0\leq\alpha\leq1$ a family of projections $E(\alpha)$
and a function $\xi(\alpha)$ with the following properties:
\begin{align*}
 \text{(i)}\quad&E(0)=0,\quad E(1)=1,
 &\text{(ii)}\quad&\alpha\leq\beta\text{ implies }E(\alpha)\leq E(\beta),\\
 \text{(iii)}\quad&\xi(0)=K,\quad\xi(1)=K^{-1},
 &\text{(iv)}\quad&\alpha\leq\beta\text{ implies }\xi(\alpha)\geq\xi(\beta),\\
 \text{(v)}\quad&D_M(E(\alpha))=\alpha,
 &\text{(vi)}\quad&E(\alpha)\text{ reduces }g,
\end{align*}
\[
 \text{(vii)}\qquad
 \alpha<\beta\text{ implies }
 \xi(\beta-0)\mathrel{\underset p{\leq}}E(\beta)-E(\alpha)
 \mathrel{\underset p{\leq}}\xi(\alpha+0).
\]
\end{statement}

Choose a sequence $\rho_1,\rho_2,\cdots$ which lies and is everywhere dense
in $0\leq\alpha\leq1$, with $\rho_1=0$, $\rho_2=1$. We will define
$E(\alpha)$ and $\xi(\alpha)$ for $\alpha=\rho_1,\rho_2,\cdots$ so that
they fulfill (i)--(vi) and
\[
 \text{(vii)$'$}\qquad
 E(\alpha)\mathrel{\underset p{\geq}}\xi(\alpha),
 \qquad 1-E(\alpha)\mathrel{\underset p{\leq}}\xi(\alpha),
\]
for all $\alpha=\rho_1,\rho_2,\cdots$.

Put first $E(0)=0$, $E(1)=1$, $\xi(0)=K$, $\xi(1)=K^{-1}$. Then
$\alpha=\rho_1,\rho_2$ are taken care of, and (i)--(vi) as well as
(vii)$'$ hold for $\rho_1,\rho_2$. Assume now that for a
$j=3,4,\cdots$ the $E(\alpha)$, $\xi(\alpha)$ for
$\alpha=\rho_1,\cdots,\rho_{j-1}$ are already defined, so that (i)--(vi),
(vii)$'$ hold for $\alpha=\rho_1,\cdots,\rho_{j-1}$; we will now define
$E(\rho_j)$, $\xi(\rho_j)$ without violating (i)--(vi), (vii)$'$.

Consider the $\rho_i$, $i=1,\cdots,j-1$ with $\rho_i\leq\rho_j$ (they
exist: $\rho_1=0\leq\rho_j$); let $\rho_{i'}$ be the greatest one. Consider
the $\rho_i$, $i=1,\cdots,j-1$, with $\rho_i\geq\rho_j$ (they exist:
$\rho_2=1\geq\rho_j$); let $\rho_{i''}$ be the smallest one. As
$\rho_{i'},\rho_{i''}\ne\rho_j$, so we have
$\rho_{i'}<\rho_j<\rho_{i''}$. So (ii) gives
$E(\rho_{i'})\leq E(\rho_{i''})$, and (v) gives
\(D_M(E(\rho_{i''})-E(\rho_{i'}))
=\rho_{i''}-\rho_{i'}>\rho_j-\rho_{i'}>0\).

Apply now \StatementRef{lemma:3.4.1}{Lemma~3.4.1} to
\(E=E(\rho_{i''})-E(\rho_{i'})\),
\(\lambda=\rho_j-\rho_{i'}\),
with $a=\xi(\rho_{i''})$, $b=\xi(\rho_{i'})$. Owing to (vii)$'$ and
\StatementRef{lemma:1.2.3}{Lemma~1.2.3} we have
$E(\rho_{i''})-E(\rho_{i'})\leq E(\rho_{i''})$ and therefore
$\mathrel{\underset p{\geq}}\xi(\rho_{i''})$, and
$E(\rho_{i''})-E(\rho_{i'})\leq1-E(\rho_{i'})$ and therefore
$\mathrel{\underset p{\leq}}\xi(\rho_{i'})$. That is,
$a\mathrel{\underset p{\leq}}E\mathrel{\underset p{\leq}}b$. An
$E_0\leq E(\rho_{i''})-E(\rho_{i'})$ and a $\xi_0$ result; put
\(E(\rho_j)=E(\rho_{i'})+E_0\) (this is a projection in $M$) and
\(\xi(\rho_j)=\xi_0\).
Let us now consider (i)--(vi), (vii)$'$ for
$\alpha=\rho_1,\cdots,\rho_{j-1},\rho_j$.

(i), (iii) are unaffected.

In (ii), (iv) the only new possibilities are
$\alpha=\rho_i\leq\rho_j=\beta$ and
$\alpha=\rho_j\leq\rho_i=\beta$, where $i=1,\cdots,j-1$. In the first case
$\rho_i\leq\rho_{i'}$, so
\(E(\alpha)=E(\rho_i)\leq E(\rho_{i'})
\leq E(\rho_{i'})+E_0=E(\rho_j)=E(\beta)\), and
\(\xi(\alpha)=\xi(\rho_i)\geq\xi(\rho_{i'})=b
\geq\xi_0=\xi(\rho_j)=\xi(\beta)\).
In the second case $\rho_i\geq\rho_{i''}$, so
\(E(\beta)=E(\rho_i)\geq E(\rho_{i''})
=E(\rho_{i'})+E(\rho_{i''})-E(\rho_{i'})
\geq E(\rho_{i'})+E_0=E(\rho_j)=E(\alpha)\), and
\(\xi(\beta)=\xi(\rho_i)\leq\xi(\rho_{i''})=a
\leq\xi_0=\xi(\rho_j)=\xi(\alpha)\).
So (ii), (iv) remain true.

In (v) we need only $D_M(E(\rho_j))=\rho_j$. Now
\(D_M(E(\rho_j))=D_M(E(\rho_{i'})+E_0)
=D_M(E(\rho_{i'}))+D_M(E_0)
=\rho_{i'}+(\rho_j-\rho_{i'})=\rho_j\).

In (vi) we need only that $E(\rho_j)$ reduces $g$. As
$E(\rho_j)=E(\rho_{i'})+E_0$, and $E(\rho_{i'})$ reduces $g$, we must only
show that $E_0$ reduces $g$ (use
\StatementRef{lemma:3.3.3}{Lemma~3.3.3}). But $E_0$ reduces $g$ with
respect to $E=E(\rho_{i''})-E(\rho_{i'})$ by
\StatementRef{lemma:3.4.1}{Lemma~3.4.1}, (iii), and
$E(\rho_{i''})-E(\rho_{i'})$ reduces $g$ because $E(\rho_{i''})$ and
$E(\rho_{i'})$ do (again use
\StatementRef{lemma:3.3.3}{Lemma~3.3.3}), therefore $E_0$ reduces $g$ by
\StatementRef{lemma:3.3.2}{Lemma~3.3.2}. So the property (vi) remains true.

In (vii)$'$ we need only
\(E(\rho_j)\mathrel{\underset p{\geq}}\xi(\rho_j)\),
\(1-E(\rho_j)\mathrel{\underset p{\leq}}\xi(\rho_j)\).
Now $E(\rho_j)=E(\rho_{i'})+E_0$, and
\(E(\rho_{i'})\mathrel{\underset p{\geq}}\xi(\rho_{i'})
\geq\xi(\rho_j)\),
\(E_0\mathrel{\underset p{\geq}}\xi_0=\xi(\rho_j)\)
by \StatementRef{lemma:3.4.1}{Lemma~3.4.1}, (iv), so
$E(\rho_j)\mathrel{\underset p{\geq}}\xi(\rho_j)$ by
\StatementRef{lemma:3.3.6}{Lemma~3.3.6}.
On the other hand
\(1-E(\rho_j)=1-E(\rho_{i''})
+\bigl((E(\rho_{i''})-E(\rho_{i'}))-E_0\bigr)\)
and by \StatementRef{lemma:3.4.1}{Lemma~3.4.1}, (iv),
\(1-E(\rho_{i''})\mathrel{\underset p{\leq}}\xi(\rho_{i''})
\leq\xi(\rho_j)\),
\((E(\rho_{i''})-E(\rho_{i'}))-E_0
=E-E_0\mathrel{\underset p{\leq}}\xi_0=\xi(\rho_j)\),
so that $1-E(\rho_j)\mathrel{\underset p{\leq}}\xi(\rho_j)$ by
\StatementRef{lemma:3.3.6}{Lemma~3.3.6}. Therefore (vii)$'$ remains true.

Thus we have verified (i)--(vi), (vii)$'$ for
$\alpha=\rho_1,\cdots,\rho_{j-1},\rho_j$. Therefore all
$\alpha=\rho_1,\rho_2,\cdots$ are taken care of, and (i)--(vi), (vii)$'$
hold for all of them. Owing to (ii) and (iv)
\(\lim_{\substack{\rho_i\to\alpha\\\rho_i\geq\alpha}}E(\rho_i)\) and
\(\lim_{\substack{\rho_i\to\alpha\\\rho_i\geq\alpha}}\xi(\rho_i)\)
exist for all $\alpha$ with $0\leq\alpha\leq1$. If $\alpha$ is equal to a
$\rho_j$, then this limit is meant to denote the value at $\rho_j$. So we
can extend the definitions of $E(\alpha)$ and $\xi(\alpha)$ to all above
$\alpha$ by defining
\[
 E(\alpha)=\lim_{\substack{\rho_i\to\alpha\\\rho_i\geq\alpha}}E(\rho_i),
 \qquad
 \xi(\alpha)=\lim_{\substack{\rho_i\to\alpha\\\rho_i\geq\alpha}}\xi(\rho_i).
\]

The statements (i), (iii) are unaffected by this extension, while (ii),
(iv)--(vi) extend by continuity to all $0\leq\alpha\leq1$.

Let us now consider (vii). Assume $\alpha<\beta$. Choose a sequence
$\rho_{i_n}$, $n=0,\mathord\pm1,\mathord\pm2,\cdots$, with
\(\alpha<\cdots<\rho_{i_{-2}}<\rho_{i_{-1}}<\rho_{i_0}
<\rho_{i_1}<\rho_{i_2}<\cdots<\beta\)
and
\(\lim_{n\to-\infty}\rho_{i_n}=\alpha\),
\(\lim_{n\to\infty}\rho_{i_n}=\beta\).
We have $\rho_{i_{-1}}>\rho_{i_{-2}}>\cdots>\alpha$, so that
\(E(\rho_{i_{-1}})\geq E(\rho_{i_{-2}})\geq\cdots\geq E(\alpha)\).
Therefore $\lim_{n\to-\infty}E(\rho_{i_n})$ exists, and is
$\geq E(\alpha)$. Now we have
\begin{align*}
 D_M\left(\lim_{n\to-\infty}E(\rho_{i_n})-E(\alpha)\right)
 &=\lim_{n\to-\infty}D_M(E(\rho_{i_n}))-D_M(E(\alpha))\\
 &=\lim_{n\to-\infty}\rho_{i_n}-\alpha=\alpha-\alpha=0,\\
 \lim_{n\to-\infty}E(\rho_{i_n})&=E(\alpha).
\end{align*}
Similarly
\[
 \lim_{n\to\infty}E(\rho_{i_n})=E(\beta).
\]
Therefore
\begin{align*}
 \sum_{m=-\infty}^{\infty}
 \bigl(E(\rho_{i_{m+1}})-E(\rho_{i_m})\bigr)
 &=\lim_{n\to\infty}\sum_{m=-n}^{n-1}
 \bigl(E(\rho_{i_{m+1}})-E(\rho_{i_m})\bigr)\\
 &=\lim_{n\to\infty}\bigl(E(\rho_{i_n})-E(\rho_{i_{-n}})\bigr)
 =E(\beta)-E(\alpha).
\end{align*}

We have
$E(\rho_{i_{m+1}})-E(\rho_{i_m})\leq E(\rho_{i_{m+1}})$, therefore
(vii)$'$ (for $\alpha=\rho_{i_{m+1}}$) and
\StatementRef{lemma:1.2.3}{Lemma~1.2.3} give
\(E(\rho_{i_{m+1}})-E(\rho_{i_m})
\mathrel{\underset p{\geq}}\xi(\rho_{i_{m+1}})
\geq\xi(\beta-0)\).
Then \StatementRef{lemma:3.3.6}{Lemma~3.3.6} gives
\(E(\beta)-E(\alpha)\mathrel{\underset p{\geq}}\xi(\beta-0)\).
On the other hand we have
$E(\rho_{i_{m+1}})-E(\rho_{i_m})\leq1-E(\rho_{i_m})$; therefore
(vii)$'$ (for $\alpha=\rho_{i_m}$) and
\StatementRef{lemma:1.2.3}{Lemma~1.2.3} give
\(E(\rho_{i_{m+1}})-E(\rho_{i_m})
\mathrel{\underset p{\leq}}\xi(\rho_{i_m})
\leq\xi(\alpha+0)\).
Then \StatementRef{lemma:3.3.6}{Lemma~3.3.6} gives
\(E(\beta)-E(\alpha)\mathrel{\underset p{\leq}}\xi(\alpha+0)\).
Thus we have proved
\[
 \xi(\beta-0)\mathrel{\underset p{\leq}}E(\beta)-E(\alpha)
 \mathrel{\underset p{\leq}}\xi(\alpha+0),
\]
that is, (vii).

We have established (i)--(vii) for all $0\leq\alpha\leq1$, and so the
proof is completed.

\begin{statement}{lemma:3.4.3}{Lemma 3.4.3}
Let $\alpha_i$, $i=0,1,\cdots,n$ be a set of numbers such that
$0=\alpha_0\leq\alpha_1\leq\cdots\leq\alpha_{n-1}\leq\alpha_n=1$
and let $\lambda_i$, $i=1,\cdots,n$ be a set of positive real numbers. Put
(with the $E(\alpha)$, $\xi(\alpha)$ of
\StatementRef{lemma:3.4.2}{Lemma~3.4.2})
\[
 g'=\sum_{i=1}^n\lambda_i(E(\alpha_i)-E(\alpha_{i-1}))g,
\]
\[
 C_1=\max_{i=1,\cdots,n}\lambda_i^2\xi(\alpha_{i-1}+0),
 \qquad
 C_2=\min_{i=1,\cdots,n}\lambda_i^2\xi(\alpha_i-0).
\]
Then for all definite $A\in M$
\[
 C_1\TrM(A)\geq(Ag',g')\geq C_2\TrM(A).
\]
\end{statement}

We have
\begin{align*}
 (Ag',g')
 &=\left(A\sum_{i=1}^n\lambda_i(E(\alpha_i)-E(\alpha_{i-1}))g,
 \sum_{j=1}^n\lambda_j(E(\alpha_j)-E(\alpha_{j-1}))g\right)\\
 &=\sum_{i=1}^n\sum_{j=1}^n\lambda_i\lambda_j
 \bigl(A(E(\alpha_i)-E(\alpha_{i-1}))g,
 (E(\alpha_j)-E(\alpha_{j-1}))g\bigr)\\
 &=\sum_{i=1}^n\sum_{j=1}^n\lambda_i\lambda_j
  \bigl((E(\alpha_j)-E(\alpha_{j-1}))A
  (E(\alpha_i)-E(\alpha_{i-1}))g,g\bigr).
\end{align*}
As every $E(\alpha_k)-E(\alpha_{k-1})$ reduces $g$ (by
\StatementRef{lemma:3.4.2}{Lemma~3.4.2}, (vi), and
\StatementRef{lemma:3.3.3}{Lemma~3.3.3}) and as
$
 (E(\alpha_j)-E(\alpha_{j-1}))\cdot
 (E(\alpha_i)-E(\alpha_{i-1}))=0
$ for $j\ne i$,
\StatementRef{lemma:3.3.4}{Lemma~3.3.4} permits us to drop all terms with
$i\ne j$ from the above sum $\sum_{i,j=1}^n$. Therefore it becomes
$
 \sum_{i=1}^n\lambda_i^2
 \bigl((E(\alpha_i)-E(\alpha_{i-1}))A
 (E(\alpha_i)-E(\alpha_{i-1}))g,g\bigr).
$

Now
$
 E(\alpha_i)-E(\alpha_{i-1})
 \mathrel{\underset p{\leq}}\xi(\alpha_{i-1}+0)
$
by \StatementRef{lemma:3.4.2}{Lemma~3.4.2}, (vii). Therefore it is evident
from \StatementRef{lemma:1.4.1}{Lemma~1.4.1} that
$\bigl((E(\alpha_i)-E(\alpha_{i-1}))A
 (E(\alpha_i)-E(\alpha_{i-1}))g,g\bigr)
 \leq\xi(\alpha_{i-1}+0)
 \TrM\bigl((E(\alpha_i)-E(\alpha_{i-1}))A
 (E(\alpha_i)-E(\alpha_{i-1}))\bigr)$. So
\begin{align*}
 (Ag',g')
 &\leq\sum_{i=1}^n\lambda_i^2\xi(\alpha_{i-1}+0)
 \TrM\bigl((E(\alpha_i)-E(\alpha_{i-1}))A
 (E(\alpha_i)-E(\alpha_{i-1}))\bigr)\\
 &\leq C_1\sum_{i=1}^n
 \TrM\bigl((E(\alpha_i)-E(\alpha_{i-1}))A
 (E(\alpha_i)-E(\alpha_{i-1}))\bigr)\\
 &\leq C_1\TrM(A)
\end{align*}
(use \StatementRef{lemma:3.3.5}{Lemma~3.3.5}). Similarly
$E(\alpha_i)-E(\alpha_{i-1})
\mathrel{\underset p{\geq}}\xi(\alpha_i-0)$
by \StatementRef{lemma:3.4.2}{Lemma~3.4.2}, (vii). Therefore
\StatementRef{lemma:1.4.1}{Lemma~1.4.1} gives
$\bigl((E(\alpha_i)-E(\alpha_{i-1}))A
 (E(\alpha_i)-E(\alpha_{i-1}))g,g\bigr)
 \geq\xi(\alpha_i-0)
 \TrM\bigl((E(\alpha_i)-E(\alpha_{i-1}))A
 (E(\alpha_i)-E(\alpha_{i-1}))\bigr)$. So
\begin{align*}
 (Ag',g')
 &\geq\sum_{i=1}^n\lambda_i^2\xi(\alpha_i-0)
 \TrM\bigl((E(\alpha_i)-E(\alpha_{i-1}))A
 (E(\alpha_i)-E(\alpha_{i-1}))\bigr)\\
 &\geq C_2\sum_{i=1}^n
 \TrM\bigl((E(\alpha_i)-E(\alpha_{i-1}))A
 (E(\alpha_i)-E(\alpha_{i-1}))\bigr)
 =C_2\TrM(A)
\end{align*}
(use \StatementRef{lemma:3.3.5}{Lemma~3.3.5}). Thus the proof is
completed.

\SectionStart{3.5}{3.5. Construction of an exact g for Tr M(A). Removal of α ≥ 1. Unrestricted weak continuity of Tr M(A). Theorems II, III, IV}3.5. We can now take the decisive step.

\begin{statement}{lemma:3.5.1}{Lemma 3.5.1}
Let $g$ and $K$ be as in \StatementRef{theorem:I}{Theorem~I}. Let $K'$ be
any number such that $1<K'\leq K$. Then there is a $g'$ which is related to
$K'$ as $g$ is to $K$ in \StatementRef{theorem:I}{Theorem~I}, and such
that
\[
 \lVert g'-g\rVert\leq(K-1)K^{1/2}.
\]
\end{statement}

Let $n$ be the smallest positive integer with $(K')^n\geq K$. Put
$\theta=K^{1/n}$, so
$1<\theta\leq K'$. Choose the $\alpha_i$, $i=0,1,\cdots,n$ with
$0=\alpha_0\leq\alpha_1\leq\cdots\leq\alpha_{n-1}\leq\alpha_n=1$, so that
$\xi(\alpha_i+0)\leq\theta^{n-2i}\leq\xi(\alpha_i-0)$.
(For $i=0$, $\theta^{n-2i}=\theta^n=K$; for $i=n$,
$\theta^{n-2i}=\theta^{-n}=K^{-1}$.) Put
$\lambda_i=\theta^{i-(n+1)/2}$. Now apply
\StatementRef{lemma:3.4.3}{Lemma~3.4.3}. We have
$
 \lambda_i^2\xi(\alpha_{i-1}+0)
 \leq\theta^{2i-(n+1)}\theta^{n-2(i-1)}
 =\theta\leq K'
$ and
$
 \lambda_i^2\xi(\alpha_i-0)
 \geq\theta^{2i-(n+1)}\cdot\theta^{n-2i}
 \geq\theta^{-1}\geq(K')^{-1},
$
so $C_1\leq K'$ and $C_2\geq(K')^{-1}$.
Therefore the $g'$ of that lemma gives for every definite $A\in M$:
$
 K'\TrM(A)\geq(Ag',g')
 \geq(K')^{-1}\TrM(A)$; that is,
$
 K'(Ag',g')\geq\TrM(A)\geq(K')^{-1}(Ag',g').
$
So it is related to $K'$ as $g$ and $K$ are in
\StatementRef{theorem:I}{Theorem~I}.

Compute now $\lVert g'-g\rVert$. The projections
$E(\alpha_i)-E(\alpha_{i-1})$, $i=1,\cdots,n$, are mutually orthogonal,
therefore the $(E(\alpha_i)-E(\alpha_{i-1}))g$, $i=1,\cdots,n$ are
mutually orthogonal too. Now we have
\begin{align*}
 \lVert g'-g\rVert^2
 &=\left\lVert
 \sum_{i=1}^n\lambda_i(E(\alpha_i)-E(\alpha_{i-1}))g
 -\sum_{i=1}^n(E(\alpha_i)-E(\alpha_{i-1}))g
 \right\rVert^2\\
 &=\left\lVert\sum_{i=1}^n(\lambda_i-1)
 (E(\alpha_i)-E(\alpha_{i-1}))g\right\rVert^2\\
 &=\sum_{i=1}^n(\lambda_i-1)^2
 \lVert(E(\alpha_i)-E(\alpha_{i-1}))g\rVert^2\\
 &\leq\max_{i=1,\cdots,n}(\lambda_i-1)^2
 \sum_{i=1}^n\lVert(E(\alpha_i)-E(\alpha_{i-1}))g\rVert^2\\
 &=\max_{i=1,\cdots,n}(\lambda_i-1)^2\lVert g\rVert^2.
\end{align*}
% Editorial note: source prints exponent (n+1)/2 here, despite the preceding
% definition of \lambda_i; retained verbatim.
Now clearly
\[
 \max_{i=1,\cdots,n}(\lambda_i-1)^2
 =(\lambda_n-1)^2
 =\bigl(\theta^{(n+1)/2}-1\bigr)^2
 \leq(\theta^n-1)^2=(K-1)^2,
\]
so that $\lVert g'-g\rVert\leq(K-1)\lVert g\rVert$.

But $A=1$ in \StatementRef{theorem:I}{Theorem~I} gives
\[
 \lVert g\rVert^2=(g,g)\leq K\TrM(1)=K,
 \qquad \lVert g\rVert\leq K^{1/2}.
\]
Therefore $\lVert g'-g\rVert\leq(K-1)K^{1/2}$. Thus the proof is
completed.

% Editorial note: source reads lowercase k in ``as g is to k''; retained
% verbatim.
Let $K_i=1+2^{-(i+2)}$, $i=0,1,2,\cdots$. We define inductively a sequence
of elements $g_i$, $i=0,1,2,\cdots$. Let $g_0$ be related to $K_0$ as $g$
is to $k$ in \StatementRef{theorem:I}{Theorem~I}. Suppose $g_i$ has been
defined and is related to $K_i$ as in \StatementRef{theorem:I}{Theorem~I}.
Then in \StatementRef{lemma:3.5.1}{Lemma~3.5.1} let
$g=g_i$, $K=K_i$, $K'=K_{i+1}$. We then define $g_{i+1}$ as $g'$. Then
$g_{i+1}$ and $K_{i+1}$ are related as $g$ and $K$ in
\StatementRef{theorem:I}{Theorem~I} and we also have
\[
 \lVert g_{i+1}-g_i\rVert
 \leq(K_i-1)K_i^{1/2}
 =\frac1{2^{i+2}}K_i^{1/2}
 \leq\frac1{2^{i+1}}.
\]

Then for $p>0$
\begin{align*}
 \lVert g_{n+p}-g_n\rVert
 &=\left\lVert\sum_{i=n}^{n+p-1}(g_{i+1}-g_i)\right\rVert
 \leq\sum_{i=n}^{n+p-1}\lVert g_{i+1}-g_i\rVert\\
 &\leq\sum_{i=n}^{n+p-1}\frac1{2^{i+1}}
 \leq\sum_{i=n}^{\infty}\frac1{2^{i+1}}=\frac1{2^n}.
\end{align*}
This implies that the $g_i$'s satisfy the Cauchy condition. Let $g$ be
their limit, and let $A$ again be definite and in $M$. Then, since $A$ is
bounded,
\[
 (Ag,g)=\lim_{i\to\infty}K_i(Ag_i,g_i)
 \geq\TrM(A)
 \geq\lim_{i\to\infty}K_i^{-1}(Ag_i,g_i)=(Ag,g).
\]
Thus, if $A\in M$ is definite, $(Ag,g)=\TrM(A)$.

For $A\in M$, Hermitian, we have as in \SectionRef{2.1}, under
\StatementRef{property:II}{Property II}, $A=A_1-A_2$, $A_1$ and $A_2$, in $M$ and positive definite.
\StatementRef{property:I}{Property I} of \SectionRef{2.1} then yields
$(Ag,g)=\TrM(A)$. If $A$ is merely in $M$, then
\StatementRef{definition:2.2.1}{Definition~2.2.1} now implies that
$(Ag,g)=\TrM(A)$. We have thus demonstrated

\begin{statement}{theorem:II}{Theorem II}
Let $M$ be a factor in case II, with $\alpha\geq1$. (In the normalization of
\SectionRef{1.1}. In the normalization of \RO{}, Theorem~VIII
this means $M,M'$ are either in case $\mathrm{II}_1,\mathrm{II}_\infty$ or
in case $\mathrm{II}_1,\mathrm{II}_1$ with $C\leq1$, standard
normalization.) Then there exists a $g\in\mathfrak H$ such that
\[
 (Ag,g)=\TrM(A).
\]
\end{statement}

We now drop the restriction that $\alpha\geq1$. Let $M,M'$ be in case
$\mathrm{II}_1,\mathrm{II}_1$, $\alpha<1$, and let $m$ be an integer such
that $m\alpha\geq1$ or in the standard normalization of \RO{}, Theorem~X,
$C/m\leq1$. Let $E_1,E_2,\cdots,E_m$ be $m$ projections each of which is
in $M$, $D_M(E_i)=1/m$, $E_iE_j=0$ if $i\ne j$. Let the range of $E_i$ be
$\mathfrak M_i$. Let us recall \RO{}, \S11.3 using $\mathfrak M_i$ for
$\mathfrak M$.

Consider the $M_{(\mathfrak M_i)}=M_i$ and
$M'_{(\mathfrak M_i)}=M_i'$ of \RORef{Definition~11.3.1}. By
\RORef{Lemma~11.3.2}, these
constitute a factorization in $\mathfrak M_i$. Now we can take for $A\in M$
and such that $E_iA=AE_i=A$,
$D_{M_i}(AE_i)=D_M(A)$,
% Editorial note: after A'\in M', the source prints AE_i (without a prime)
% on the left while retaining A' on the right; retained verbatim.
and for $A'\in M'$,
$D_{M_i'}(AE_i)=D_{M'}(A')$,
$D_{M_i}$ and $D_{M_i'}$ are dimension functions for $M_i$ and $M_i'$ 
respectively (\RORef{Lemma~11.3.6}).

The ranges of $D_{M_{(\mathfrak M)}}$ and $D_{M'_{(\mathfrak M)}}$ are by
\RORef{Lemma~11.3.7}, the intervals $(0,1/m)$, $(0,\alpha)$ respectively. So we can
obtain the standard normalization of \RO{}, Theorem~X by multiplying by
$m$ and $1/\alpha$ respectively. But in this latter normalization, by
\RORef{Lemma~11.4.3},
$C=\bigl(D_M(E_i)/D_{M'}(1)\bigr)C_0=C_0/m\leq1$,
where $D_M$ and $D_{M'}$ refer to the standard normalization for $M$ and
$M'$ as in \RO{}, Theorem~X.

By \StatementRef{theorem:II}{Theorem~II} above this implies that there is a
$g_i^0\in\mathfrak M_i$ such that, if $A_0\in M_i$, then
$Tr_{M_i}(A_0)=(A_0g_i^0,g_i^0)$. 
But now if $A\in M$ is such that $E_iA=AE_i=A$, let $A_0=AE_i$. From the
above it will be seen that
$\TrM(A)=\frac1mTr_{M_i}(A_0)=\frac1m(A_0g_i^0,g_i^0)
=\frac1m(Ag_i^0,g_i^0)=(Ag_i,g_i)$, if $g_i=(1/m)^{1/2}g_i^0$.

% Editorial note: source again cites the apparently nonexistent Lemma 3.2.6;
% the printed citation is retained, while its link points to Lemma 3.3.5,
% whose trace decomposition is used here.
Now if $A\in M$, we have by \StatementRef{lemma:3.3.5}{Lemma~3.2.6}
\begin{align*}
 \TrM(A)
 &=\sum_{i=1}^m\TrM(E_iAE_i)
 =\sum_{i=1}^m(E_iAE_ig_i,g_i)\\
 &=\sum_{i=1}^m(AE_ig_i,E_ig_i)
 =\sum_{i=1}^m(Ag_i,g_i).
\end{align*}

If $M,M'$ is in case $\mathrm I_m,\mathrm I_{m'}$, $m<\infty$, we may
proceed as follows. By \RO{}, Lemma~8.6.1,
$\mathfrak H=E_m\otimes\mathfrak H_2$, where $E_m$ is an $m$-dimensional
euclidean space. Let $\phi_1,\cdots,\phi_m$ be a complete orthonormal set in
$E_m$ and $f\in\mathfrak H_2$ with $\lVert f\rVert=1$. Then if we let
$g_i=\phi_i\otimes f$, we easily verify that the above formula for
$\TrM$ holds. Thus we have

\begin{statement}{theorem:III}{Theorem III}
If $M$ is a factor in a finite case, then there exists a finite number of
elements $g_i\in\mathfrak H$ such that
\[
 \TrM(A)=\sum_{i=1}^m(Ag_i,g_i).
\]
\end{statement}

Suppose that $A\in M$ is held fixed. Then consider the weak neighborhood
$U=U(A;g_1,\cdots,g_m;g_1,\cdots,g_m;\epsilon/m)$.
$X\in U$ implies $|((X-A)g_i,g_i)|<\epsilon/m$
for $i=1,\cdots,m$. Hence if $X$ is also in $M$,
\StatementRef{theorem:III}{Theorem~III} implies
$|\TrM(X)-\TrM(A)|<\epsilon$. This proves

\begin{statement}{theorem:IV}{Theorem IV}
If $M$ is a factor in a finite case, $\TrM(A)$ is weakly
continuous.
\end{statement}

\chapterhead{IV}{The isomorphism of $M$, $M'$ and $\mathfrak H$
  (for $\alpha=1$)}{Chapter IV. The isomorphism of M, M', and H (for α = 1)}

\SectionStart{4.1}{4.1. The isomorphism of H and Q(M) with the help of a u.d.f.}4.1. We assume that $M,M'$ is in case
$\mathrm{II}_1,\mathrm{II}_1$ and $\alpha=1$ or what is the same thing that
$C=1$ and we have the standard normalization. We know by the discussion of
\RO{}, \S\S11.3 and 11.4 how all other cases II can be reduced to this case.

As $\alpha=1$, \StatementRef{theorem:II}{Theorem~II} holds. We define

\begin{statement}{definition:4.1.1}{Definition 4.1.1}
A $g$ which satisfies \StatementRef{theorem:II}{Theorem~II} is uniformly
distributed with respect to $M$; abbreviated $g$ u.d.r. $M$.
\end{statement}

\begin{statement}{lemma:4.1.1}{Lemma 4.1.1}
If $g$ is u.d.r. $M$, then (i) $\lVert g\rVert=1$,
(ii) $E_g^{M'}=1$, (iii) $E_g^M=1$ (cf. \RO{}, Definition~5.1.1).
\end{statement}

To prove (i), in \StatementRef{theorem:II}{Theorem~II} take $A$ as 1. Then
$1=\TrM(1)=(g,g)=\lVert g\rVert^2$.

To prove (ii), in \StatementRef{theorem:II}{Theorem~II} take $A$ as
$E_g^{M'}$. Then
\[
 1=\lVert g\rVert^2=(E_g^{M'}g,g)
 =\TrM(E_g^{M'})=D_M(E_g^{M'})
\]
or $D_M(E_g^{M'})=1$ which in our normalization implies $E_g^{M'}=1$.



% PDF p. 29; journal p. 236
Consider (iii). Since we have the standard normalization and \(C=1\),
$D_{M'}(E_g^M)=D_M(E_g^{M'})=1$.
This implies \(E_g^M=1\).

\begin{statement}{def:4.1.2}{Definition 4.1.2}
Let \(g\) be \(\in\cH\). Let \(\Q_g(M)\) consist of all operators \(Z\in U(M)\),
i.e., \(Z\) is linear closed \(\eta M\) and with a dense domain
(cf. \RO, Definition 4.2.1), for which \(Zg\) exists.
\end{statement}

Now consider any \(f\in\cH\) such that \(E_f^{M'}=E_f^M=1\). By \RO,
Lemma 9.2.1, if \(h\in\cH\), then \(h=XYf\), where \(X\) and \(Y\) are
% Editorial note: source has no comma after case II_1 here; retained verbatim.
\(\in U(M)\). But if \(M\) is in case \(\mathrm{II}_1\) by \RO, Theorem XV,
\(Z=[XY]\) is \(\in U(M)\), and, since \(Z\supseteq XY\), \(Zf\) exists, and
we also have \(h=XYf=Zf\); i.e., every \(h\in\cH\) is in the form \(Zf\),
where \(Z\) is \(\in U(M)\).

Furthermore this \(Z\) is uniquely determined by \(h\). For suppose \(Z_1\)
and \(Z_2\in U(M)\) are such that \(Z_1f=Z_2f\) or
\((Z_1-Z_2)f=0\). Then \([Z_1-Z_2]\) is by \RO, Theorem XV, \(\in U(M)\)
and since \([Z_1-Z_2]\supseteq Z_1-Z_2\), \([Z_1-Z_2]f\) exists and is zero.
Let \(A'\) be \(\in M'\). Then since
\(A'[Z_1-Z_2]\subseteq[Z_1-Z_2]A'\),
$[Z_1-Z_2]A'f=A'[Z_1-Z_2]f=0$.
% Editorial note: source prints A'\in M here (without the prime on M); retained verbatim.
Thus \([Z_1-Z_2]\) is zero on the set of \(A'f\), \(A'\in M\). But since
\(E_f^{M'}=1\), this latter set is everywhere dense in \(\cH\). Thus
\([Z_1-Z_2]\) is zero on a dense set and since it is closed it must be zero.
\RO, Theorem XV, then yields
\[
 Z_1=[Z_1-0]=[Z_1-[Z_1-Z_2]]
     =[[Z_1-Z_1]+Z_2]=[0+Z_2]=Z_2.
\]
So we have proved

\begin{statement}{lem:4.1.2}{Lemma 4.1.2}
If \(f\) is such that \(E_f^{M'}=E_f^M=1\) (in particular, if \(f\) is
u.d.r.\ \(M\)), then \(h=Zf\) defines a one-to-one correspondence of
\(\cH\) and the set of all \(Z\in U(M)\) for which \(Zf\) exists
(i.e., between \(\cH\) and \(\Q_f(M)\)).
\end{statement}

Since our assumptions on \(f\) are symmetric in \(M\) and \(M'\), we also
have

\begin{statement}{lem:4.1.3}{Lemma 4.1.3}
Let \(f\) be as in \StatementRef{lem:4.1.2}{Lemma 4.1.2}. Then \(h=Z'f\)
defines a one-to-one correspondence of \(\cH\) and the set of all
\(Z'\in U(M')\) for which \(Z'f\) exists (i.e., between \(\cH\) and
\(\Q_f(M')\)).
\end{statement}

Thus for \(h\in\cH\) we have three correspondences defined by the equations
\(Zf=h=Z'f\), if \(f\) is u.d.r.\ \(M\): (1) the correspondence
\(\mathfrak I_M\) of \(\cH\) and \(\Q_f(M)\), (2) the correspondence
\(\mathfrak I_{M'}\) of \(\cH\) and \(\Q_f(M')\), and (3) the correspondence
\(\mathfrak I_{M,M'}\) of \(\Q_f(M)\) and \(\Q_f(M')\).

\SectionStart{4.2}{4.2. The correspondences of H, Q(M), and Q(M'). Criteria for u.d. Equivalence of u.d. for M and for M'. Isomorphism properties of these correspondences. Theorems V, VI}4.2.\quad
We now investigate these three correspondences. We know that they are
one-to-one. They are obviously linear and of course along with \(\cH\) the
sets \(\Q_f(M)\) and \(\Q_f(M')\) are linear. But inasmuch as
\(\mathfrak I_{M,M'}\) is a correspondence of operators, we would like to
know the algebraic properties of this correspondence as well as certain
properties of \(\Q_f(M)\) and \(\Q_f(M')\). In particular we would like to
know:
\begin{enumerate}[label=(\roman*),leftmargin=2.6em]
\item To what extent can the operation \([XY]\) be performed in
      \(\Q_f(M)\) (or \(\Q_f(M')\))?
\item To what extent can the operation \(X^*\) be performed in
      \(\Q_f(M)\) (or \(\Q_f(M')\))?
\item Is the property of being bounded invariant under
      \(\mathfrak I_{M,M'}\); i.e., since \(M\subset\Q_f(M)\),
      \(M'\subset\Q_f(M')\), is \(M\) mapped on \(M'\) by
      \(\mathfrak I_{M,M'}\)?
\item Are the properties of being Hermitian, definite, or unitary invariant
      under \(\mathfrak I_{M,M'}\)?
\end{enumerate}

% PDF p. 30; journal p. 237
\begin{statement}{lem:4.2.1}{Lemma 4.2.1}
If \(f\) is u.d.r.\ \(M\), then every \(Uf\), \(U\in M\) and unitary, is
u.d.r.\ \(M\) also.
\end{statement}

If \(A\in M\), then
\[
 (AUf,Uf)=(U^{-1}AUf,f)=\Tr_M(U^{-1}AU)=\Tr_M(A).
\]

\begin{statement}{lem:4.2.2}{Lemma 4.2.2}
If \(f\) and \(g\) are each u.d.r.\ \(M\), then there exists a
\(U'\in M'\) and unitary with \(g=U'f\).
\end{statement}

If \(h=Af\), \(A\in M\), define \(h^*\) as \(Ag\). The correspondence
\(h\rightleftarrows h^*\) is linear, and, since
\begin{align*}
 \lVert h^*\rVert^2
 &=\lVert Ag\rVert^2=(Ag,Ag)=(A^*Ag,g)=\Tr_M(A^*A)=(A^*Af,f)\\
 &=(Af,Af)=\lVert Af\rVert^2=\lVert h\rVert^2,
\end{align*}
it is isometric. Thus the equation \(Wh=h^*\) defines a linear isometric and
hence one-valued operator \(W\). Since \(E_f^M=E_g^M=1\), both domain and
range of \(W\) are everywhere dense and hence its closure \(\widetilde W\)
is unitary.

Now if \(h\) is in the domain of \(W\), \(h=Af\), \(A\in M\). Then, for every
\(B\in M\),
$BWh=Bh^*=BAg=WBAf=WBh$
or \(BW\subseteq WB\). This implies that \([BW]\subseteq[WB]\). But
\(B\widetilde W=\widetilde{BW}\subseteq[BW]\), while since \(B\) is bounded
\([WB]=\widetilde WB\). Hence \(B\widetilde W\subseteq\widetilde WB\). If in
particular \(B\) is unitary, \RO, Lemma 4.2.2 now yields that
\(\widetilde W\) is \(\eta M'\) and \RO, Lemma 4.2.1 implies that
\(\widetilde W\) is \(\in M'\).

Letting \(A=1\) in the definition of \(W\), we get
\(g=Wf=\widetilde Wf\).

\begin{statement}{lem:4.2.3}{Lemma 4.2.3}
If \(f\) and \(g\) are u.d.r.\ \(M\), then there exists a \(U\in M\) and
unitary such that \(g=Uf\).
\end{statement}

By \StatementRef{lem:4.1.2}{Lemma 4.1.2}, \(g=Zf\) where \(Z\) is
\(\in U(M)\). Using the canonical decomposition for \(Z\), \(Z=BW\), where
\(B\) is self-adjoint and definite and \(W\) is partially isometric and as
indicated in the proof of \StatementRef{lemma:1.4.3}{Lemma 1.4.3} may be taken as unitary; that is,
\(g=BWf\), where \(W\) is unitary and \(B\) definite and self-adjoint. We
must show that \(B=1\).

By \StatementRef{lem:4.2.1}{Lemma 4.2.1}, \(f_0=Wf\) is u.d.r.\ \(M\).
Hence we must show that if \(f_0\) and \(g\) are u.d.r.\ \(M\) and
\(g=Bf_0\), where \(B\) is definite and self-adjoint, then \(B=1\).
% PDF p. 31; journal p. 238
Let \(E(\lambda)\) be the resolution of the identity corresponding to \(B\).
If \(E(\lambda)=0\) for \(\lambda<1\) and \(E(\lambda)=1\) for
\(\lambda>1\), then \(B=1\). Thus if \(B\) is not \(1\), either there exists
a \(\lambda_1<1\) for which \(E(\lambda_1)\ne0\) or a \(\lambda_2>1\) for
which \(E(\lambda_2)\ne1\). In the first case we have
$
 \lVert E(\lambda_1)f_0\rVert^2
  =(E(\lambda_1)f_0,E(\lambda_1)f_0)
    =(E(\lambda_1)f_0,f_0)=\Tr_M(E(\lambda_1))
  =(E(\lambda_1)g,g)=\lVert E(\lambda_1)g\rVert^2
    =\lVert E(\lambda_1)Bf_0\rVert^2
    =\lVert BE(\lambda_1)f_0\rVert^2
  \leq \lambda_1^2\lVert E(\lambda_1)f_0\rVert^2
$ or
$(1-\lambda_1^2)\lVert E(\lambda_1)f_0\rVert^2\leq0$.
This implies
$
 0=\lVert E(\lambda_1)f_0\rVert^2
   =(E(\lambda_1)f_0,f_0)=\Tr_M(E(\lambda_1))
   =D_M(E(\lambda_1))
$
or \(E(\lambda_1)=0\), a contradiction. In the second case a similar type
of calculation yields that
$
 \lVert(1-E(\lambda_2))f_0\rVert^2
 =\lVert B(1-E(\lambda_2))f_0\rVert^2
 \geq\lambda_2^2\lVert(1-E(\lambda_2))f_0\rVert^2,
$
which again implies that \(1-E(\lambda_2)=0\), a contradiction. Thus \(B\)
must be \(1\) and our lemma is proved.

\begin{statement}{lem:4.2.4}{Lemma 4.2.4}
If \(f\) is u.d.r.\ \(M\), and \(U'\in M'\) and unitary, then \(U'f\) is
u.d.r.\ \(M\).
\end{statement}

If \(A\) is \(\in M\),
\[
 (AU'f,U'f)=(U'^{-1}AU'f,f)=(AU'^{-1}U'f,f)=(Af,f)=\Tr_M(A).
\]

\begin{statement}{lem:4.2.5}{Lemma 4.2.5}
Unitary operators correspond to unitary operators under
\(\mathfrak I_{M,M'}\); i.e., if \(f\) is u.d.r.\ \(M\), \(U\in M\),
\(U'\in M'\) and \(Uf=U'f\), then if either \(U\) or \(U'\) is unitary the
other is too.
\end{statement}

Let \(U\in M\) be unitary, then by \StatementRef{lem:4.2.1}{Lemma 4.2.1}
\(Uf\) is u.d.r.\ \(M\). Then \StatementRef{lem:4.2.2}{Lemma 4.2.2} yields
% Editorial note: source reads U'\in M here; retained verbatim.
that there is a unitary \(U'\in M\), such that \(U'f=Uf\). Thus
\(U\sim U'\) under \(\mathfrak I_{M,M'}\). Similarly
\StatementRef{lem:4.2.4}{Lemmas 4.2.4} and
\StatementRef{lem:4.2.3}{4.2.3} yield that to every unitary
% Editorial note: source reads U'\in M here; retained verbatim.
\(U'\in M\),
there exists a unitary \(U\in M\) such that \(U'f=Uf\).

\begin{statement}{lem:4.2.6}{Lemma 4.2.6}
\(f\) u.d.r.\ \(M\) is equivalent to \(f\) u.d.r.\ \(M'\).
\end{statement}

% Editorial note: source reads A'\in M here; retained verbatim.
Let \(f\) be u.d.r.\ \(M\) and \(A'\in M\). Then consider
$T_0(A')=(A'f,f)$.
It is linear in \(A'\). By \StatementRef{lemma:4.1.1}{Lemma 4.1.1}, (i) \(T_0(1)=1\). If \(A'\) is
% Editorial note: source abruptly uses f_0 here instead of f; retained verbatim.
definite, then \(T_0(A')=(A'f_0,f_0)\geq0\). Furthermore we have
$
 T_0(A'^*)=(A'^*f,f)=(f,A'f)
 =\overline{(A'f,f)}
% Editorial note: source reads \overline{T_0(A)} here; retained verbatim.
 =\overline{T_0(A)}.
$
Therefore \(T_0(A')\) satisfies (i)--(iv) and (v) of
\StatementRef{property:III}{Property III} with
\(M'\) instead of \(M\).

We would like to verify next (vi)\('\) of the same property for \(T_0(A')\).
By \StatementRef{lem:4.2.5}{Lemma 4.2.5} we have for any unitary
% Editorial note: source reads U'\in M here; retained verbatim.
\(U'\in M\)
\begin{align*}
 T_0(U'^{-1}A'U')
  &=(U'^{-1}A'U'f,f)=(A'U'f,U'f)=(A'Uf,Uf)\\
  &=(U^{-1}A'Uf,f)=(A'(U^{-1}U)f,f)=(A'f,f)=T_0(A').
\end{align*}
Thus (vi)\('\) holds. \StatementRef{property:IV}{Property IV} of
\SectionRef{2.2} now implies that
\(T_0(A')=\Tr_{M'}(A')\) and hence \(f\) is u.d.r.\ \(M'\).

Replacing \(M\) by \(M'\) in the above argument yields the converse.

We note that the above argument also proves

\begin{statement}{lem:4.2.7}{Lemma 4.2.7}
If \(f\in\cH\) is such that (i) \(\lVert f\rVert=1\) and (ii) if to every
unitary \(U\in M\) there exists a unitary \(U'\in M'\) such that \(Uf=U'f\),
then \(f\) is u.d.r.\ \(M\).
\end{statement}

% PDF p. 32; journal p. 239
\begin{statement}{thm:V}{Theorem V}
Let \(f\in\cH\) be such that \(\lVert f\rVert=1\) and let \(M,M'\) be in
case \(\mathrm{II}_1,\mathrm{II}_1\) with \(C=1\). Then the following
statements are equivalent:
\begin{enumerate}[label={},leftmargin=2.4em]
\item[\((\alpha)\)] \(f\) is u.d.r.\ \(M\).
\item[\((\beta)\)] \(f\) is u.d.r.\ \(M'\).
\item[\((\gamma)\)] If \(U\) is unitary and \(\in M\), then there exists a unitary
      \(U'\in M'\) such that \(Uf=U'f\).
\item[\((\delta)\)] If \(U'\) is unitary and \(\in M'\), then there exists a unitary
      \(U\in M\) such that \(U'f=Uf\).
\end{enumerate}

Furthermore either \((\alpha)\), \((\beta)\), \((\gamma)\), or
\((\delta)\) implies \(E_f^M=E_f^{M'}=1\).
\end{statement}

\((\alpha)\) is equivalent to \((\beta)\) by
\StatementRef{lem:4.2.6}{Lemma 4.2.6}. \((\alpha)\) and \((\gamma)\) are
equivalent by \StatementRef{lem:4.2.5}{Lemmas 4.2.5} and
\StatementRef{lem:4.2.7}{4.2.7}. Replacing \(M\) by \(M'\), the last
mentioned lemmas also show \((\beta)\) and \((\delta)\) to be equivalent.
This and \StatementRef{lemma:4.1.1}{Lemma 4.1.1} prove the last statement
of the theorem.

In view of \StatementRef{thm:V}{Theorem V}, we will say that \(f\) is
uniformly distributed (abbreviated u.d.) if \(f\) is u.d.r.\ \(M\).

\begin{statement}{cor:V}{Corollary}
If \(f\) is u.d., \(\mathfrak I_{M,M'}\) maps
\(M(\subseteq\Q_f(M))\) on \(M'(\subseteq\Q_f(M'))\); that is, the property
of being bounded is invariant under \(\mathfrak I_{M,M'}\).
\end{statement}

By \StatementRef{thm:V}{Theorem V}, the property of being unitary is
invariant under \(\mathfrak I_{M,M'}\). This and the linearity of
\(\mathfrak I_{M,M'}\) imply together that the property of having the form
\(A=\sum_{\nu=1}^{n}a_\nu U_\nu\)
\((n=1,2,\cdots;\ a_\nu\) complex, \(U_\nu\) unitary and in \(M\) or \(M'\)
along with \(A)\) is invariant under \(\mathfrak I_{M,M'}\).

We will show that \(A\in\Q_f(M)\) (or \(\Q_f(M')\)) is bounded if and only
if it is in this form. Now if \(A\) is in this form, it is obviously bounded
so we must only show that every bounded \(A\) is in this form. Since
\(A=A_1+iA_2\), \(A_1\) and \(A_2\) Hermitian, it will be sufficient if this
is shown for Hermitian \(A\)'s. Since if \(A\) is in this form, \(cA\) is
also, we may assume that the bound of \(A\) is \(\leq1\). Then \(1-A^2\)
is definite and \((1-A^2)^{1/2}\) exists. Letting
\(U=A+i(1-A^2)^{1/2}\),
then \(U^*=A-i(1-A^2)^{1/2}\) and \(U\) is unitary. Furthermore
\(A=\frac12(U+U^*)\).

Now since the form \(\sum_{\nu=1}^{n}a_\nu U_\nu\) is invariant under
\(\mathfrak I_{M,M'}\), boundedness must be also.

\begin{statement}{thm:VI}{Theorem VI}
Assume that \(f_0\) is u.d. Then \(\mathfrak I_{M,M'}\) is an
anti-isomorphism of \(M\) and \(M'\). That is, if \(A\sim A'\),
\(B\sim B'\), then (i) \(\alpha A\sim\alpha A'\); (ii)
\(A^*\sim A'^*\); (iii) \(A+B\sim A'+B'\); (iv) \(AB\sim B'A'\);
(v) if \(A\) (or \(A'\)) is Hermitian, then \(A'\) (or \(A\)) is also.
\end{statement}

(i) and (iii) are obvious.

Consider (iv). We have
\(ABf_0=AB'f_0=B'Af_0=B'A'f_0\),
so that \(AB\sim B'A'\). To prove (v), let \(U\) be unitary,
\(U\sim U'\). Then \(U'\) is unitary also by
\StatementRef{thm:V}{Theorem V}. Similarly inasmuch as \(U^{-1}\) is
unitary, if \(U^{-1}\sim V'\), \(V'\) is unitary. As
% PDF p. 33; journal p. 240
\(UU^{-1}=U^{-1}U=1\) by (iv) we have \(U'V'=V'U'=1\), and hence
\(V'=(U')^{-1}\). So \(U^{-1}\sim(U')^{-1}\) and
\(\alpha(U+U^{-1})\sim\alpha(U'+(U')^{-1})\)
by (iii). If \(\alpha\) is \(>0\), the right side is Hermitian and the left
side is any Hermitian element of \(M\) (cf. the proof of the
\StatementRef{cor:V}{corollary} of
\StatementRef{thm:V}{Theorem V}). This proves (v).

We now prove (ii). If \(A\) is \(\in M\), then
\(A=B+iC\), \(A^*=B-iC\), \(B,C\) are Hermitian and \(\in M\). If
\(A\sim A'\), \(B\sim B'\), \(C\sim C'\), then \(B',C'\) are Hermitian and
\(A\sim B'+iC'=A'\), \(A^*\sim B'-iC'=A'^*\),
proving (ii).

\begin{statement}{cor:VI}{Corollary}
The properties of being Hermitian, being definite, being a projection, as
well as the numerical quantities
\(\opnorm{A}\) (the bound of \(A\)) and
\(\Tr_M(A)\) (or \(\Tr_{M'}(A')\)), all within \(M\) (or \(M'\)), are
invariant under \(\mathfrak I_{M,M'}\).
\end{statement}

\(A\) Hermitian means \(A=A^*\); \(A\) a projection means
\(A=A^*=A^2\); \(A\) definite means that there exists an Hermitian \(B\)
with \(A=B^2\). So all these properties are invariant under
\(\mathfrak I_{M,M'}\). \(1\) is invariant
\((1\in M\) and \(1\in M'\), \(1f_0=1f_0)\).
\(\opnorm{A}\) is the smallest \(\alpha\geq0\) such
that \(\alpha^2\cdot1-A^*A\) is definite, so
\(\opnorm{A}\) is invariant. If \(A\sim A'\), then
\(Af_0=A'f_0\) and
% Editorial note: source reads \Tr_M(A') here; retained verbatim.
\(\Tr_M(A)=(Af_0,f_0)=(A'f_0,f_0)=\Tr_M(A')\)
by \StatementRef{thm:V}{Theorem V}.

\SectionStart{4.3}{4.3. Analysis of Q(M), Q(M'); their Hilbert-space character. Theorems VII, VIII, IX}4.3.\quad
In what follows \(f_0\) will always be assumed to be u.d. The discussion
which follows could be based on an extension of the notion of \(\Tr_M(A)\)
(and of \(\Tr_{M'}(A')\)) to unbounded \(A\,\eta M\) (or \(A'\,\eta M'\))
but we prefer an approach which avoids this.

\begin{statement}{thm:VII}{Theorem VII}
The sets \(\Q_{f_0}(M)\) and \(\Q_{f_0}(M')\) are independent of the choice
% Editorial note: source reads "choice of of"; retained verbatim.
of of u.d.\ \(f_0\). We will therefore denote them by \(\Q(M)\) and
\(\Q(M')\), respectively. Furthermore the values of
\(\lVert Zf_0\rVert\), \((Xf_0,Yf_0)\), where \(X,Y,Z\in\Q(M)\), or
\(\Q(M')\), are independent of the choice of the u.d.\ \(f_0\).
\end{statement}

Owing to the symmetry between \(M\) and \(M'\) it suffices to prove the
statements concerning \(M\); i.e., let \(f_0\) and \(g_0\) be u.d., then for
\(A\in U(M)\), \(Af_0\) is defined if and only if \(Ag_0\) is defined and
\(\lVert Af_0\rVert=\lVert Ag_0\rVert\). Owing to the symmetry in \(f_0\)
and \(g_0\) it will be sufficient to show that \(Ag_0\) is defined if
\(Af_0\) is and \(\lVert Af_0\rVert=\lVert Ag_0\rVert\).

By \StatementRef{lem:4.2.2}{Lemma 4.2.2}, there exists a unitary
\(U'\in M'\) such that \(g_0=U'f_0\). Since \(A\) is \(\in U(M)\),
\(A=U'^{-1}AU'\). Hence since \(Af_0\) exists
\(Af_0=U'^{-1}AU'f_0=U'^{-1}Ag_0\)
exists. This implies that \(Ag_0\) exists and since \(Ag_0=U'Af_0\),
\(\lVert Ag_0\rVert=\lVert Af_0\rVert\). Also if
\(B\in\Q_{f_0}(M)\), then \(Bg_0=U'Bf_0\) and
\((Ag_0,Bg_0)=(U'Af_0,U'Bf_0)=(Af_0,Bf_0)\).

So we must have these mappings:
\[
 \tag{I}\EquationTarget{map:I}
 \mathfrak I_M:\cH\sim\Q(M);\qquad
 \mathfrak I_{M'}:\cH\sim\Q(M');\qquad
 \mathfrak I_{M,M'}:\Q(M)\sim\Q(M').
\]
By the \StatementRef{cor:V}{corollary} to \StatementRef{thm:V}{Theorem V},
\(\mathfrak I_{M,M'}\) maps \(M\) on \(M'\). So \(\mathfrak I_M\) maps
\(M\) and \(\mathfrak I_{M'}\) maps \(M'\) on the same subset
\(\mathfrak A\) of \(\cH\) and we have the further mappings
\[
 \tag{II}\EquationTarget{map:II}
 \mathfrak I_M:\mathfrak A\sim M;\qquad
 \mathfrak I_{M'}:\mathfrak A\sim M';\qquad
 \mathfrak I_{M,M'}:M\sim M'.
\]
\EquationRef{map:II}{(II)} is part of \EquationRef{map:I}{(I)}.

% PDF p. 34; journal p. 241
We now prove

\begin{statement}{lem:4.3.1}{Lemma 4.3.1}
\(\mathfrak A\) is a linear set, dense in \(\cH\).
\end{statement}

Inasmuch as \(M\) is linear, \(\mathfrak A\) is too.

Consider an \(f\in\cH\), \(f=Zf_0\), \(Z\in\Q(M)\). Now \(Z=BU\), \(B\)
self-adjoint and definite, \(B\in U(M)\), \(U\) unitary and \(\in M\).
Write \(B\) in its spectral form
\[
 B=\int_0^\infty \lambda\,dE(\lambda),\qquad E(\lambda)\in M.
\]
Then we have
\[
 [E(\mu)B]=\int_0^\mu \lambda\,dE(\lambda);
\]
and therefore (i) \([E(\mu)B]\in M\) and (ii) if \(Bg\) is defined then we
have
\(\mathop{\mathrm{strong}\ \lim}_{\mu\to\infty}\allowbreak[E(\mu)B]g\allowbreak=Bg\).
Thus \([E(\mu)B]U\in M\) and
\(\mathop{\mathrm{strong}\ \lim}_{\mu\to\infty}\allowbreak[E(\mu)B]Uf_0
\allowbreak=BUf_0\allowbreak=Zf_0\allowbreak=f\).
Since \([E(\mu)B]Uf_0\in\mathfrak A\), \(f\) is a condensation point of
\(\mathfrak A\). As this is true for any \(f\in\cH\), \(\mathfrak A\) is
dense in \(\cH\).

\begin{statement}{def:4.3.1}{Definition 4.3.1}
If \(f=Xf_0=X'f_0\), \(g=Yf_0=Y'f_0\), \(X\) and \(Y\in\Q(M)\), \(X'\)
and \(Y'\in\Q(M')\), let
\([\![X]\!]=[\![X']\!]=\lVert f\rVert\),
\(\langle\!\langle X,Y\rangle\!\rangle
=\langle\!\langle X',Y'\rangle\!\rangle=(f,g)\).
\end{statement}

\StatementRef{thm:VII}{Theorem VII} shows that the values of
\([\![X]\!]\), \([\![X']\!]\),
\(\langle\!\langle X',Y'\rangle\!\rangle\),
\(\langle\!\langle X,Y\rangle\!\rangle\) are independent of the choice of
the u.d.\ \(f_0\).

\begin{statement}{lem:4.3.2}{Lemma 4.3.2}
If \(X\) and \(Y\) are \(\in M\), then
\([\![X]\!]=(\Tr_M(X^*X))^{1/2}=(\Tr_M(XX^*))^{1/2}\),
\(\langle\!\langle X,Y\rangle\!\rangle
=\Tr_M(Y^*X)=\Tr_M(XY^*)\).
\end{statement}

As
\([\![X]\!]=\bigl(\langle\!\langle X,X\rangle\!\rangle\bigr)^{1/2}\),
the second statement implies the first. Now clearly
\[
\langle\!\langle X,Y\rangle\!\rangle
 =(Xf_0,Yf_0)=(Y^*Xf_0,f_0)=\Tr_M(Y^*X)
\]
% Editorial note: source writes f here instead of the standing f_0; retained verbatim.
since \(f\) is u.d. We also have \(\Tr_M(Y^*X)=\Tr_M(XY^*)\) by
\SectionRef{2.2}, \StatementRef{property:III}{Property III}, (vi). So
\(\langle\!\langle X,Y\rangle\!\rangle
=\Tr_M(Y^*X)=\Tr_M(XY^*)\).

\begin{statement}{lem:4.3.3}{Lemma 4.3.3}
If \(X'\), \(Y'\in M'\), then
\([\![X']\!]=(\Tr_{M'}(X'^*X'))^{1/2}
=(\Tr_{M'}(X'X'^*))^{1/2}\),
\(\langle\!\langle X',Y'\rangle\!\rangle=\Tr_{M'}(Y'^*X')
% Editorial note: source reads \Tr_M(X'Y'^*) here; retained verbatim.
=\Tr_M(X'Y'^*)\).
\end{statement}

Replace \(M\) by \(M'\) in the proof of
\StatementRef{lem:4.3.2}{Lemma 4.3.2}.

\begin{statement}{thm:VIII}{Theorem VIII}
If we use the definitions
\begin{align*}
 \langle\!\langle X,Y\rangle\!\rangle
   &=\Tr_M(Y^*X)=\Tr_M(XY^*),\\
 [\![X]\!]
   &=\bigl(\langle\!\langle X,X\rangle\!\rangle\bigr)^{1/2}
     =(\Tr_M(X^*X))^{1/2}=(\Tr_M(XX^*))^{1/2},
\end{align*}
then \(M\) is an incomplete Hilbert space. Its completion in the usual
(Cantor) way gives a Hilbert space \(\widetilde M\) which may be identified
with \(\Q(M)\).
\end{statement}

% PDF p. 35; journal p. 242
This is clear by the isomorphisms \EquationRef{map:I}{(I)} and
\EquationRef{map:II}{(II)} and \StatementRef{lem:4.3.1}{Lemmas 4.3.1} and
\StatementRef{lem:4.3.2}{4.3.2}. \(M\) and \(\widetilde M\), that is,
\(\Q(M)\), are isomorphic to \(\mathfrak A\) and \(\cH\) respectively,
using the isomorphism \(\mathfrak I_M\).

The corresponding facts hold for \(M'\); their formulation is obvious.

\begin{statement}{thm:IX}{Theorem IX}
\(X\in\Q(M)\) implies \(X^*\in\Q(M)\) and
\([\![X]\!]=[\![X^*]\!]\).
\end{statement}

Assume \(X\in\Q(M)\), i.e., \(X\in U(M)\); \(Xf_0\) is defined. Write
\(X=BU\), \(B\) self-adjoint and \(\eta M\), \(U\) unitary and \(\in M\).
Thus \(BUf_0\) is defined. Now \(Uf_0\) is u.d.\
(\StatementRef{lem:4.2.1}{Lemma 4.2.1}), hence the existence of \(BUf_0\)
implies \(B\in\Q(M)\) (\StatementRef{thm:VII}{Theorem VII}). This in turn
implies the existence of \(Bf_0\) and with it
\(X^*f_0=U^{-1}Bf_0\). So \(X^*\in\Q(M)\).

As \(f_0\) and \(Uf_0\) are both u.d.\
(\StatementRef{lem:4.2.1}{Lemma 4.2.1}) we have by
\StatementRef{def:4.3.1}{Definition 4.3.1},
$
 [\![X]\!]=\lVert Xf_0\rVert=\lVert BUf_0\rVert
 =[\![B]\!]=\lVert Bf_0\rVert=\lVert U^{-1}Bf_0\rVert
 =\lVert X^*f_0\rVert=[\![X^*]\!].
$

\begin{statement}{cor:IX}{Corollary}
\(X\sim X^*\) is an involutory conjugate anti-isomorphism of \(\Q(M)\)
and isometric.
\end{statement}

\(X\sim X^*\) maps \(\Q(M)\) on part of itself by
\StatementRef{thm:IX}{Theorem IX}. This and \(X^{**}=X\) show that
\(X\sim X^*\) is a one-to-one and involutory mapping of \(\Q(M)\) on
itself. It is isometric by \StatementRef{thm:IX}{Theorem IX};
\([\![X]\!]=[\![X^*]\!]\) and it is obviously a conjugate
anti-isomorphism.

The corresponding facts hold for \(M'\); the formulation is obvious.

\SectionStart{4.4}{4.4. Algebraic properties of Q(M). Roles of M and M' in it. Properties I⁰-V⁰. Theorem X}4.4.\quad
We next discuss the algebra of \(\Q(M)\).

\begin{statement}{prop:I0}{Property \(\mathrm{I}^{0}\)}
\(A,B\in\Q(M)\) imply \(\alpha A,A^*,[A+B]\in\Q(M)\). If also either
\(A\) or \(B\) is \(\in M\), then \([AB]\) is \(\in\Q(M)\).
\end{statement}

\(\alpha A\), \([A+B]\) are \(\in\Q(M)\) since \(\Q(M)\) is linear
(along with \(\cH\)). \(A^*\) is \(\in\Q(M)\) by
\StatementRef{thm:IX}{Theorem IX}.

Assume now that \(A\in M\), \(B\in\Q(M)\). Then \(Bf_0\) is defined and so
is
\(ABf_0=[AB]f_0\).
Thus \([AB]\in\Q(M)\). If \(A\in\Q(M)\), \(B\in M\), then \(B^*\in M\),
\(A^*\in\Q(M)\), so \([B^*A^*]\in\Q(M)\). Since
\([B^*A^*]^*=[AB]\) (\RO, Theorem XV), \([AB]\) is \(\in\Q(M)\).

\begin{statement}{prop:II0}{Property \(\mathrm{II}^{0}\)}
(i) \([\![\alpha A]\!]=|\alpha|\cdot[\![A]\!]\); (ii)
\([\![A^*]\!]=[\![A]\!]\); (iii)
\([\![A+B]\!]\leq[\![A]\!]+[\![B]\!]\); (iv)
% Editorial note: source omits a second comparison sign after "and";
% the printed formula is retained verbatim.
\([\![[AB]]\!]\leq\opnorm{A}\cdot[\![B]\!]\) and
\([\![A]\!]\cdot\opnorm{B}\).
\end{statement}

(i) and (iii) hold because \(\Q(M)\) is isomorphic to \(\cH\). (ii) holds
by \StatementRef{thm:IX}{Theorem IX}.

Consider (iv). We have
\[
 [\![[AB]]\!]=\lVert[AB]f_0\rVert=\lVert ABf_0\rVert
 \leq\opnorm{A}\cdot\lVert Bf_0\rVert
 =\opnorm{A}\cdot[\![B]\!]
\]
% PDF p. 36; journal p. 243
proving the first inequality. The second follows from this, by using (ii)
\[
 [\![[AB]]\!]=[\![[AB]^*]\!]=[\![[B^*A^*]]\!]
 \leq\opnorm{B^*}\cdot[\![A^*]\!]
 =\opnorm{B}\cdot[\![A]\!].
\]

We have \(\cH\sim\Q(M)\). Neither \(\cH\) nor \(\Q(M)\) depends on the
choice of the u.d.\ \(f_0\) but \(\mathfrak I_M\) does. This dependence is
as follows.

\begin{statement}{prop:III0}{Property \(\mathrm{III}^{0}\)}
Let us replace the u.d.\ \(f_0\) by a u.d.\ \(g_0\) and let
\(\widetilde{\mathfrak I}_M\) denote the resulting correspondence. Let
\(U\) be unitary and \(\in M\) and such that \(Uf_0=g_0\)
% Editorial note: the source cites Lemma 4.2.1 here, though Lemma 4.2.3
% supplies the stated unitary for the given f_0 and g_0. The printed citation
% is retained, while its link points to Lemma 4.2.3.
(cf. \StatementRef{lem:4.2.3}{Lemma 4.2.1}). Then if \(X=YU\), \(X\) and
\(Y\in\Q(M)\), we have
\[
 \mathfrak I_M:f\sim X;\qquad
 \widetilde{\mathfrak I}_M:f\sim Y
\]
if \(f=Xf_0\).
\end{statement}

This is clear since \(f=Xf_0=YUf_0=Yg_0\).

We can now determine another notion which does not depend on the choice of
the u.d.\ \(f_0\).

\begin{statement}{prop:IV0}{Property \(\mathrm{IV}^{0}\)}
The linear set \(\mathfrak A\) (cf. \SectionRef{4.3}, the correspondences
\EquationRef{map:II}{(II)}) does not depend on the choice of the u.d.\
\(f_0\).
\end{statement}

This follows from the definition of \(\mathfrak A\) and the fact that if
\(X=YU\), \(U\) unitary and either \(X\) or \(Y\) is bounded, then both
\(X\) and \(Y\) are bounded.

Now \(\Q(M)\sim\Q(M')\) under \(\mathfrak I_{M,M'}\) and while neither
\(\Q(M)\) nor \(\Q(M')\) depends on the choice of the u.d.\ \(f_0\) yet
\(\mathfrak I_{M,M'}\) does. We now obtain this dependence.

\begin{statement}{prop:V0}{Property \(\mathrm{V}^{0}\)}
Let the u.d.\ \(f_0\) be replaced by the u.d.\ \(g_0\) and let the resulting
correspondence between \(\Q(M)\) and \(\Q(M')\) be denoted by
\(\widetilde{\mathfrak I}_{M,M'}\). Then if \(U\in M\) is unitary and such
that \(g_0=Uf_0\), and \(X=U^{-1}YU\), \(X\) and \(Y\in\Q(M)\), then
\[
 \mathfrak I_{M,M'}:X'\sim X;\qquad
 \widetilde{\mathfrak I}_{M,M'}:X'\sim Y
\]
if \(X'f_0=Xf_0\), \(X'\in\Q(M')\).
\end{statement}

Let \(X'f_0=Xf_0\), then
\(X'g_0=X'Uf_0=UX'f_0=UXf_0=UXU^{-1}g_0\).

\begin{statement}{def:4.4.1}{Definition 4.4.1}
The isomorphism \(\cH\sim\Q(M)\) makes correspond to every operator \(P\)
in \(\cH\) an operator \(P^0\) in \(\Q(M)\).
\end{statement}

Observe that inasmuch as the elements of \(\Q(M)\) are operators in
\(\cH\), \(P^0\) is an operator on the operators of \(\cH\).

\begin{statement}{thm:X}{Theorem X}
\((i)\) If \(A\in M\), then
\[
 A^0Z=[AZ];
\]

\((ii)\) if \(A'\in M'\) and \(A'\sim A\in M\) under
\(\mathfrak I_{M,M'}\), then
\[
 % Editorial note: source omits the argument Z on the left; the proof below
 % prints A'{}^0Z; the theorem is retained verbatim.
 A'{}^0=[ZA].
\]
\end{statement}

% PDF p. 37; journal p. 244
\(P^0\) is defined as follows. If \(g=Xf_0\), \(h=Pg\), \(h=Yf_0\),
\(X\) and \(Y\in\Q(M)\), then \(P^0X=Y\). Thus to prove (i), we have
\((A^0Z)f_0=A(Zf_0)=[AZ]f_0\)
and to show (ii)
\((A'{}^0Z)f_0=A'Zf_0=ZA'f_0=ZAf_0=[ZA]f_0\).

\begin{statement}{cor:X:1}{Corollary 1}
Let \(M^0\) be the set of all operators \(L_A\) in \(\Q(M)\), \(A\in M\),
where \(L_AZ=[AZ]\), and \(M_0\) the set of all operators \(R_A\) in
\(\Q(M)\), \(A\in M\), where \(R_AZ=[ZA]\). Then \(\mathfrak I_M\) carries
\(M\) into \(M^0\) and \(M'\) into \(M_0\).
\end{statement}

This is obvious by \StatementRef{thm:X}{Theorem X}.

\begin{statement}{cor:X:2}{Corollary 2}
\(M^0,M_0\) are rings and \(M^{0\prime}=M_0\) (all in \(\Q(M)\)). They are
factors of class \((\mathrm{II}_1,\mathrm{II}_1)\) with \(C=1\).
\end{statement}

Since \(\mathfrak I_M\) is a spatial isomorphism of \(\cH\) and \(\Q(M)\)
which takes \(M,M'\) into \(M^0,M_0\) respectively, these properties hold.

We have now characterized \(\cH,M,M'\) by the situation in
\(\Q(M),M^0,M_0\) to which they are spatially isomorphic. The spatial
isomorphism is \(\mathfrak I_M\) which depends on an arbitrary u.d.\
\(f_0\), while \(\Q(M),M^0,M_0\) themselves do not. All the influence of
the choice of \(f_0\) is however a further transformation \(R_U\),
\(X=YU\) (\(U\in M\), \(U\) unitary so \(R_U\in M_0\)).
\(\mathfrak I_M\) always carries the totality of u.d.\ elements \(g_0\) of
\(\cH\) into the totality of unitary elements \(U\in\Q(M)\), \(f_0\) being
that element which corresponds to \(1\).

\SectionStart{4.5}{4.5. Implication of a spatial isomorphism by an algebraic ring-isomorphism. Theorem XI}4.5.\quad
The following important isomorphism theorem can now be proved.

\begin{statement}{thm:XI}{Theorem XI}
Let \(\cH_1\) and \(\cH_2\) be two Hilbert spaces and \(M_1,M_1'\) and
\(M_2,M_2'\) respectively be factor pairs in them, both of class
\((\mathrm{II}_1,\mathrm{II}_1)\) and with \(C=1\) in the standard
normalization. Then an algebraic ring isomorphism of \(M_1\) and \(M_2\)
is a necessary and sufficient condition for the existence of a spatial
isomorphism of \(\cH_1\) and \(\cH_2\) which takes \(M_1,M_1'\) into
\(M_2,M_2'\).
\end{statement}

The necessity is obvious and so we prove the sufficiency. Let
\(\mathfrak F\) be the algebraic ring isomorphism of \(M_1\) and \(M_2\).
By \SectionRef{2.2}, \StatementRef{property:IV}{Property IV},
\(\Tr_{M_1}(A_1)=\Tr_{M_2}(A_2)\) if \(A_1\sim A_2\) under
\(\mathfrak F\). So
\([\![A_1]\!]=[\![A_2]\!]\),
\(\langle\!\langle A_1,B_1\rangle\!\rangle
=\langle\!\langle A_2,B_2\rangle\!\rangle\),
% Editorial note: source reads B_2\sim B_2 here; retained verbatim.
if \(A_1\sim A_2\), \(B_2\sim B_2\) under \(\mathfrak F\). Thus
\(\mathfrak F\) is an isometric mapping of \(M_1\) on \(M_2\) and therefore
% Editorial note: source reads Q(M) on Q(M') here; retained verbatim.
extends by continuity in a unique way to a linear and isometric mapping of
\(\Q(M)\) on \(\Q(M')\) which we call \(\mathfrak F\) again.
\(A_1\sim A_2\), \(B_1\sim B_2\) under \(\mathfrak F\) imply
\(A_1B_1\sim A_2B_2\) under \(\mathfrak F\) if \(A_1,B_1\in M_1\) (and
thus \(A_2,B_2\in M_2\)). By continuity this will even hold, if one of
\(A_1,B_1\) is in \(M_1\) and the other merely in \(\Q(M_1)\) (and one of
\(A_2,B_2\) in \(M_2\), and the other merely in \(\Q(M_2)\)). So
\(\mathfrak F\) is a spatial isomorphism of \(\Q(M_1)\) and \(\Q(M_2)\)
which carries \(M_1^0,M_{1,0}\) into \(M_2^0,M_{2,0}\) (by
\StatementRef{cor:X:1}{Corollary 1} to
% Editorial note: the preceding source citation reads Theorem XI, although
% the numbered corollary appears under Theorem X. The printed citation is
% retained, while its link points to Theorem X.
\StatementRef{thm:X}{Theorem XI}). Now (by the same
\StatementRef{cor:X:1}{corollary})
\(\mathfrak I_{M_2}\mathfrak F\mathfrak I_{M_1}^{-1}\) is a spatial
isomorphism of \(\cH_1\) and \(\cH_2\) which carries \(M_1,M_1'\) into
\(M_2,M_2'\).

\SectionStart{4.6}{4.6. Description of further results}\StatementRef{thm:XI}{Theorem XI} could be extended to cover cases
\((\mathrm{II}_1,\mathrm{II}_1)\), where
\(C\mathrel{\substack{\geq\\<}}1\) in the standard
normalization as well as other combinations of \(\mathrm{II}_1\) and
\(\mathrm{II}_\infty\). In all
% PDF p. 38; journal p. 245
cases a reduction of the spatial-isomorphism questions of \(\cH_{\!1}\) and
% Editorial note: source says algebraic-isomorphism questions of H_1 and H_2,
% rather than of M_1 and M_2; retained verbatim.
\(\cH_{\!2}\) to algebraic-isomorphism questions of \(\cH_{\!1}\) and
\(\cH_{\!2}\)
(plus the behavior of \(C\)) results. We will discuss these questions in
detail in later publications.

\appendixhead

\SectionStart{A.1}{1.}1.\quad
Consider a ring \(M\) of class \(\mathrm{II}_1\) in \(\cH\). Then \(M'\) is
of class \(\mathrm{II}_1\) or \(\mathrm{II}_\infty\). Normalize \(D_M\) and
\(D_{M'}\), so as to have \(D_M(\cH)=1\) and \(C=1\), so
\(D_{M'}(\cH)=\alpha\) (cf. \SectionRef{1.1}). By choosing an
\(n=1,2,\cdots\) with \(1/n\leq\alpha\) and then an
\(\mathfrak M'\,\eta\,M'\) with \(D_{M'}(\mathfrak M')=1/n\), apply \RO,
\S11.3, to form \(M_{(\mathfrak M')}\), \(M'_{(\mathfrak M')}\) in
\(\mathfrak M'\). Then \(M_{(\mathfrak M')}\) is
algebraically-ring-isomorphic to \(M\), and
\(D_{M_{(\mathfrak M')}}(\mathfrak M')=D_M(\cH)=1\)
(cf. \RO, Lemmas 11.3.3 and 11.3.6). So if we are interested in the
algebraical properties of \(M\) only, we may assume without any loss in
generality that \(M,M'\) are in case \(\mathrm{II}_1,\mathrm{II}_1\) and
that \(\alpha=1/n\), \(n=1,2,\cdots\). Or, in standard normalization
\(C=n\), \(n=1,2,\cdots\) (cf. \RO, Theorem X).

Now form the direct product of \(\cH\) with an \(n\)-dimensional Euclidean
space \(E_n\oplus\cH\) as in \SectionRef{2.1}, for the case \(C>1\), and
consider \(M^{(2)}\) in \(E_n\oplus\cH\). The argument used in
\SectionRef{2.1} shows that
% Editorial note: source repeats M^{(2)\prime} on both sides; retained verbatim.
\(M^{(2)\prime}=R(N^{(1)},M^{(2)\prime})\) and \(C^{(2)}=C/n=1\)
% Editorial note: source prints ``is.ring'' with the intervening period; retained verbatim.
and \(M^{(2)}\) is.ring isomorphic to \(M\).

So we have for \(M^{(2)},M^{(2)\prime}\) in the standard normalization
\(C=1\). If we are therefore interested in the algebraical properties of
\(M\) only, we may even assume without any loss in generality that \(C=1\)
in standard normalization; that is, \(\alpha=1\) in the normalization of
\SectionRef{1.1}. We will assume this in what follows.

It may be noticed concerning subrings that the metric of \(\Q(M)\)
(cf. \StatementRef{def:4.3.1}{Definition} \StatementRef{def:4.3.1}{4.3.1}) is determined by the
algebra of \(M\) (cf. \SectionRef{2.2}). A subring demands closure in the
weak operator topology, but it will be shown elsewhere that for subrings of
\(M\) weak, relative closure in \(M\) for the \([\![X-Y]\!]\) metric is
equivalent to closure in the weak topology.

\SectionStart{A.2}{2.}2.\quad
Under these conditions there is an analogy between \(M\) and the matrices
of a Euclidean space \(E_n\), described by an interesting
Lebesgue-Stieltjes-Radon measure in the plane.

We can use the results of \hyperlink{sec:4.2}{\(\S\S\)4.2}\hyperlink{sec:4.4}{--4.4}, and thus we may form the
Hilbert-space \(\Q(M)\) which is isomorphic to \(\cH\), and in which
\(M,M'\) are located by \StatementRef{thm:X}{Theorem X}.

Let \(E(\lambda)\), \(a<\lambda<b\), be a resolution of unity, all
\(E(\lambda)\in M\). Put
\(A=\int_{-\infty}^{\infty}\lambda\,dE(\lambda)\)
(in the well known symbolic sense). For any Borel-set of real numbers \(S\)
form
\[
 e_S(\lambda)=
 \begin{cases}
  1&\text{for }\lambda\in S,\\
  0&\text{for }\lambda\notin S,
 \end{cases}
\]
and
\[
 E(S)=e_S(A)=\int_{-\infty}^{\infty}e_S(\lambda)\,dE(\lambda)
             =\int_S dE(\lambda)
\]
% PDF p. 39; journal p. 246
(symbolically). As
$\overline{e_S(\lambda)}=e_S(\lambda)=(e_S(\lambda))^2$,
so \(e_S(A)=e_S(A)^*=(e_S(A))^2\); that is, \(E(S)=e_S(A)\in M\) is a
projection. (Cf. also Maeda, \textit{Journal of Science, Hiroshima
University}, Ser.\ A, vol.\ 4 (1934), pp.\ 57--91.)

Consider now point sets in the \(\lambda,\mu\)-plane \(P\) and in
particular sets of the form \(S_1\otimes S_2\):
\[
 (\lambda,\mu)\in S_1\otimes S_2
 \quad\hbox{means}\quad
 \lambda\in S_1,\ \mu\in S_2,
\]
\(S_1,S_2\) being two Borel-sets of real numbers. Define for any \(X\in M\)
(or even \(X\in\Q(M)\)),
\[
 X_{S_1\times S_2}=E(S_1)XE(S_2)
\]
and
\begin{align*}
 \omega(X;S_1\otimes S_2)
  &=[\![X_{S_1\times S_2}]\!]^2
    =\Tr\bigl((X_{S_1\times S_2})^*X_{S_1\times S_2}\bigr)\\
  &=\Tr\bigl(E(S_2)X^*E(S_1)\cdot E(S_1)XE(S_2)\bigr)\\
  &=\Tr\bigl(E(S_2)X^*\cdot E(S_1)XE(S_2)\bigr)\\
  &=\Tr\bigl(E(S_1)XE(S_2)\cdot E(S_2)X^*\bigr)\\
  &=\Tr\bigl(E(S_1)XE(S_2)X^*\bigr).
\end{align*}

The first expression for \(\omega(X,S_1\otimes S_2)\) shows, that it is
always \(\geq0\), the last one, that it is a totally additive set-function
of \(S_1\) and of \(S_2\).

Denote the set of all \(\lambda\) by \(\Lambda\), then the above facts imply:
\begin{align*}
 \omega(X;S_1\otimes S_2)
 &\leq\omega(X;S_1\otimes S_2)
       +\omega(X;(\Lambda-S_1)\otimes S_2)\\
 &=\omega(X;\Lambda\otimes S_2)\\
 &\leq\omega(X;\Lambda\otimes S_2)
       +\omega(X;\Lambda\otimes(\Lambda-S_2))\\
 &=\omega(X;\Lambda\otimes\Lambda)
   =[\![X_{\Lambda\times\Lambda}]\!]^2\\
 &=[\![1\cdot X\cdot1]\!]^2=[\![X]\!]^2.
\end{align*}

So we have

\begin{statement}{lem:A}{Lemma A}
\((i)\) \(\omega(X;S_1\otimes S_2)\) is defined for all linear Borel-sets
\(S_1,S_2\). \((ii)\) It is totally additive in \(S_1\) as well as in
\(S_2\). \((iii)\) We have always
\[
 0\leq\omega(X;S_1\otimes S_2)\leq[\![X]\!]^2=\Tr(X^*X).
\]
\end{statement}

\SectionStart{A.3}{3.}3.\quad
We can use \StatementRef{lem:A}{Lemma A} to define a
Lebesgue-Stieltjes-Radon measure \(\mu(X;T)\) for all plane Borel-sets
\(T(\subset P)\) with the help of \(\omega(X;S_1\otimes S_2)\).

\begin{statement}{def:A}{Definition A}
If \(T\) is a plane Borel-set (\(T\subset P\)), then consider all sequences
of linear Borel-set \(S_1^{(i)},S_2^{(i)}\), \(i=1,2,\cdots\) (all
\(S_1^{(i)},S_2^{(i)}\subset\Lambda\)) for which
\[
 \tag{*}\EquationTarget{app:eq:cover}
 T\subset\sum_{i=1}^{\infty}S_1^{(i)}\otimes S_2^{(i)}.
\]
For every such sequence form the (numerical) sum
\[
 \tag{**}\EquationTarget{app:eq:sum}
 \sum_{i=1}^{\infty}\omega(X;S_1^{(i)}\otimes S_2^{(i)}).
\]
Denote the g.l.b.\ of all numbers \EquationRef{app:eq:sum}{(**)} by
\(\mu(X;T)\).
\end{statement}

% PDF p. 40; journal p. 247
We prove now

\begin{statement}{lem:B}{Lemma B}
\((i)\) \(\mu(X;T)\) is defined for all plane Borel-sets \(T\subset P\).

\((ii)\) It is totally additive in \(T\).

\((iii)\) We have always
\[
 0\leq\mu(X;T)\leq[\![X]\!]^2=\Tr(X^*X).
\]

\((iv)\) In particular
\[
 \mu(X;S_1\otimes S_2)=\omega(X;S_1\otimes S_2)
\]
and especially
\[
 \mu(X,P)=[\![X]\!]^2=\Tr(X^*X).
\]
\end{statement}

(i) is obvious. To prove (ii), observe first that our
\StatementRef{def:A}{Definition A} makes
\[
 \mu(X;T_1+T_2+\cdots)
 \leq\mu(X;T_1)+\mu(X;T_2)+\cdots
\]
obvious, so we need to prove
\[
 \mu(X;T_1+T_2+\cdots)
 \geq\mu(X;T_1)+\mu(X;T_2)+\cdots
\]
only when \(T_i\cdot T_j=0\) for \(i\ne j\).

Even \(\mu(X;T_1+T_2)\geq\mu(X;T_1)+\mu(X;T_2)\) suffices. Then
finite induction gives
\(\mu(X;T_1+T_2+\cdots+T_n)
\geq\mu(X;T_1)+\mu(X;T_2)+\cdots+\mu(X;T_n)\), and so
\(\mu(X;T_1+T_2+\cdots)
\geq\mu(X;T_1+\cdots+T_n)
\geq\mu(X;T_1)+\cdots+\mu(X;T_n)\), and as this holds for all
\(n=1,2,\cdots\), it implies
\(\mu(X;T_1+T_2+\cdots)
\geq\mu(X;T_1)+\mu(X;T_2)+\cdots\).

We will prove
\[
 \tag{\S}\EquationTarget{app:eq:caratheodory}
 \mu(X;T)=\mu(X;TU)+\mu(X;T-TU)
\]
for all \(T,U\) this gives our above inequality if \(T=T_1+T_2\),
\(U=T_1\). Call a \(U\), for which
\EquationRef{app:eq:caratheodory}{(\S)} holds for all plane Borel-sets
\(T\), following Carath\'eodory (cf. \ROcite{4}, pp.\ 246--252), measurable.
We must show that all \(U\) are measurable.

If \(U=S_1\otimes S_2\), then
\(T\subset\sum_{i=1}^{\infty}S_1^{(i)}\otimes S_2^{(i)}\) implies
\(TU\subset\sum_{i=1}^{\infty}
S_1^{(i)}S_1\otimes S_2^{(i)}S_2\),
\(T-TU\subset\sum_{i=1}^{\infty}
S_1^{(i)}S_1\otimes\bigl(S_2^{(i)}-S_2^{(i)}S_2\bigr)
+\sum_{i=1}^{\infty}
\bigl(S_1^{(i)}-S_1^{(i)}S_1\bigr)\otimes S_2^{(i)}\).
Thus our \StatementRef{def:A}{Definition A} gives immediately
\EquationRef{app:eq:caratheodory}{(\S)} with \(\geq\) in it, and as
\(\leq\) is obvious (cf. above) it proves
\EquationRef{app:eq:caratheodory}{(\S)}. So all \(U=S_1\otimes S_2\) are
measurable, and therefore in particular all plane intervals
\((=\text{rectangles})\).

% PDF p. 41; journal p. 248
Now the measurable sets \(U\) form a Borel-ring (cf., for instance,
\ROcite{4}, loc.\ cit.). These considerations apply literally to the
present case. Thus every plane Borel-set \(U\) is measurable. This completes
the proof.

We now prove (iii). \(\mu(X;T)\geq0\) because all expressions
\EquationRef{app:eq:sum}{(**)} in \StatementRef{def:A}{Definition A} are
\(\geq0\). The relation
\(\mu(X;T)\leq[\![X]\!]^2=\Tr(X^*X)\) results by putting
\(S_1^{(1)}=S_2^{(1)}=\Lambda\) and all other
\(S_1^{(i)}=S_2^{(i)}=0\).

To prove (iv), put \(S_1^{(1)}=S_1\), \(S_2^{(1)}=S_2\) and all other
\(S_1^{(i)}=S_2^{(i)}=0\). This gives
\(\mu(X;S_1\otimes S_2)\leq\omega(X;S_1\otimes S_2)\). So we must
prove \(\geq\) only; that is,
\[
 S_1\otimes S_2
 \subset\sum_{i=1}^{\infty}S_1^{(i)}\otimes S_2^{(i)}
 \quad\hbox{implies}\quad
 \omega(X;S_1\otimes S_2)
 \leq\sum_{i=1}^{\infty}
 \omega(X;S_1^{(i)}\otimes S_2^{(i)}).
\]

Considering the properties of \(\omega(X;S_1\otimes S_2)\) given in
\StatementRef{lem:A}{Lemma A}, this implication follows literally as in
the paper of
\href{https://doi.org/10.4064/fm-23-1-237-278}{\L omnicki and Ulam}
(\textit{Fundamenta Mathematicae}, vol.\ 23 (1934), pp.\ 237--278; cf.\
also the lecture notes of the second-named author for the year 1934--1935).

So we have proved
\(\mu(X;S_1\otimes S_2)=\omega(X;S_1\otimes S_2)\). Put now
\(S_1=S_2=\Lambda\); then
\[
 \mu(X;P)=\mu(X;\Lambda\otimes\Lambda)
 =\omega(X;\Lambda\otimes\Lambda)
 =[\![X]\!]^2=\Tr(X^*X).
\]
results.

\SectionStart{A.4}{4.}4.\quad
The plane measure \(\mu(X;T)\) may be used to \lq\lq locate\rq\rq\ the
\lq\lq position\rq\rq\ of \(X\) in the \(\lambda,\mu\)-plane \(P\). It
gives the entire plane a total \lq\lq measure\rq\rq\ or
\lq\lq weight\rq\rq\ \([\![X]\!]^2\) (which is \(>0\) if \(X\ne0\)), and
any part \(T\) of it correspondingly a \(\mu(X;T)\) (\(\geq0\)). It plays
the same role, as the sum of the absolute-value-squares of all
matrix-elements in a certain area \(T\) in the matrix-scheme (which may be
looked at as a region in the plane) for the matrices of a
finite-dimensional Euclidean space \(E_n\) (that is, \(M\) in a case
\((\mathrm{I}_n)\), \(n=1,2,\cdots\)).

We will analyse it more thoroughly in subsequent publications.

\begin{center}
\textsc{Princeton University and The Institute for Advanced Study,}\\
\textsc{Princeton, N.\ J.}
\end{center}



\end{document}
